% Options for packages loaded elsewhere
\PassOptionsToPackage{unicode}{hyperref}
\PassOptionsToPackage{hyphens}{url}
\documentclass[
  12pt,
]{article}
\usepackage{xcolor}
\usepackage[margin=1in]{geometry}
\usepackage{amsmath,amssymb}
\setcounter{secnumdepth}{5}
\usepackage{iftex}
\ifPDFTeX
  \usepackage[T1]{fontenc}
  \usepackage[utf8]{inputenc}
  \usepackage{textcomp} % provide euro and other symbols
\else % if luatex or xetex
  \usepackage{unicode-math} % this also loads fontspec
  \defaultfontfeatures{Scale=MatchLowercase}
  \defaultfontfeatures[\rmfamily]{Ligatures=TeX,Scale=1}
\fi
\usepackage{lmodern}
\ifPDFTeX\else
  % xetex/luatex font selection
  \setmainfont[]{Arial}
\fi
% Use upquote if available, for straight quotes in verbatim environments
\IfFileExists{upquote.sty}{\usepackage{upquote}}{}
\IfFileExists{microtype.sty}{% use microtype if available
  \usepackage[]{microtype}
  \UseMicrotypeSet[protrusion]{basicmath} % disable protrusion for tt fonts
}{}
\makeatletter
\@ifundefined{KOMAClassName}{% if non-KOMA class
  \IfFileExists{parskip.sty}{%
    \usepackage{parskip}
  }{% else
    \setlength{\parindent}{0pt}
    \setlength{\parskip}{6pt plus 2pt minus 1pt}}
}{% if KOMA class
  \KOMAoptions{parskip=half}}
\makeatother
\usepackage{longtable,booktabs,array}
\newcounter{none} % for unnumbered tables
\usepackage{calc} % for calculating minipage widths
% Correct order of tables after \paragraph or \subparagraph
\usepackage{etoolbox}
\makeatletter
\patchcmd\longtable{\par}{\if@noskipsec\mbox{}\fi\par}{}{}
\makeatother
% Allow footnotes in longtable head/foot
\IfFileExists{footnotehyper.sty}{\usepackage{footnotehyper}}{\usepackage{footnote}}
\makesavenoteenv{longtable}
\setlength{\emergencystretch}{3em} % prevent overfull lines
\providecommand{\tightlist}{%
  \setlength{\itemsep}{0pt}\setlength{\parskip}{0pt}}
\usepackage{bookmark}
\IfFileExists{xurl.sty}{\usepackage{xurl}}{} % add URL line breaks if available
\urlstyle{same}
% fallback for those not using the hyperref driver hyperxmp:
\makeatletter
\@ifundefined{xmpquote}{\newcommand{\xmpquote}[1]{#1}}{}
\makeatother
\hypersetup{
  hidelinks,
  pdfcreator={LaTeX via pandoc}}

\author{}
\date{}

\begin{document}

{
\setcounter{tocdepth}{3}
\tableofcontents
}
\section{CONCEPTION D’UN SYSTEME D’INTELLIGENCE ARTIFICIELLE POUR LA
DETECTION DE LA FRAUDE BANCAIRE ET MOBILE MONEY AU TOGO : CAS DE SUNU
BANK}\label{conception-dun-systeme-dintelligence-artificielle-pour-la-detection-de-la-fraude-bancaire-et-mobile-money-au-togo-cas-de-sunu-bank}

Rapport en vue de l’obtention du diplôme de :

Bachelor en Intelligence Artificielle \& Big Data

Présenté par :

JOHNSON Nancy Flora Ekua Bentsiwa

Encadreur Académique : Encadreur Professionnel : Monsieur AHOULOUMI
Essowaba Monsieur TCHAKALA Fissale

LOME, TOGO Année académique : 2025-2026

\begin{center}\rule{0.5\linewidth}{0.5pt}\end{center}

\begin{quote}
\textbf{Note de révision — Version 2 (juillet 2026)}

Ce rapport a été examiné dans le cadre de la révision de fin de cycle du
Bachelor en Intelligence Artificielle \& Big Data (session juillet
2026).

\textbf{Décision :} Rapport accepté sous réserve des corrections
suivantes :

{\def\LTcaptype{none} % do not increment counter
\begin{longtable}[]{@{}
  >{\raggedright\arraybackslash}p{(\linewidth - 4\tabcolsep) * \real{0.1000}}
  >{\raggedright\arraybackslash}p{(\linewidth - 4\tabcolsep) * \real{0.3750}}
  >{\raggedright\arraybackslash}p{(\linewidth - 4\tabcolsep) * \real{0.5250}}@{}}
\toprule\noalign{}
\begin{minipage}[b]{\linewidth}\raggedright
N°
\end{minipage} & \begin{minipage}[b]{\linewidth}\raggedright
Point soulevé
\end{minipage} & \begin{minipage}[b]{\linewidth}\raggedright
Correction apportée
\end{minipage} \\
\midrule\noalign{}
\endhead
\bottomrule\noalign{}
\endlastfoot
1 & III.2.1 et III.4.1 — La mention « split 80/20 » est ambiguë (k-fold
ou holdout ?) & Remplacée par « holdout simple 80/20 (entraînement) / 20
\% (test) » dans les deux sections \\
2 & Tableau III.2 — Écart de 5 382 entre le total de la matrice (112
726) et la taille de l’échantillon de test (118 108) & Note de bas de
tableau explicitant l’écart (suppression des lignes NaN après
prédiction) \\
3 & III.4.1 — Contrainte de précision non documentée dans le code &
Ajout d’un bloc justificatif : contrainte \texttt{Précision\ ≥\ 0,15}
codée en dur dans le déclencheur d’alerte (filtre post-détection) \\
4 & Réduction du nombre de tableaux demandée & Tableau I.2 (comparaison
algorithmes) → texte ; Tableau II.2 (caractéristiques datasets) → texte
inline ; Tableau II.3 (métriques) → supprimé, formules intégrées au
texte \\
5 & Approfondissement homogène des sections & Ajout de justifications
techniques (paramètres SMOTE, écart avec le Deep Learning, arbitrage
coût/bénéfice des faux positifs) et d’une Annexe D (questions anticipées
de soutenance) \\
\end{longtable}
}

Les tableaux IV.2 et IV.3 mentionnés dans la demande de révision ne
figurent pas dans cette version du document, le chapitre IV n’étant pas
encore rédigé.

\begin{center}\rule{0.5\linewidth}{0.5pt}\end{center}
\end{quote}

\subsection{DEDICACE}\label{dedicace}

A mes parents, pour leur soutien indéfectible et leurs sacrifices tout
au long de mon parcours académique. Leur confiance et leurs
encouragements constants ont été le moteur essentiel dans l’achèvement
de mes études.

A ma famille, pour sa patience et son affection.

\begin{center}\rule{0.5\linewidth}{0.5pt}\end{center}

\subsection{REMERCIEMENTS}\label{remerciements}

Aucun travail de recherche ne saurait être l’œuvre d’une seule personne.
Sa réalisation résulte toujours du concours de plusieurs acteurs. C’est
pourquoi nous tenons à exprimer notre profonde gratitude à toutes celles
et ceux qui ont contribué à l’aboutissement de ce mémoire.

Nos remerciements s’adressent en premier lieu à notre directeur de
mémoire, Monsieur AHOULOUMI Essowaba, pour son encadrement de qualité et
son accompagnement tout au long de ce travail. Sa disponibilité, ses
orientations méthodologiques et sa rigueur scientifique ont été
déterminants dans l’aboutissement de cette recherche.

Nous exprimons également notre reconnaissance à notre encadreur
professionnel, Monsieur TCHAKALA Fissale, pour son accompagnement
technique et ses conseils avisés.

Nous tenons à exprimer notre gratitude au corps professoral et
administratif du Collège de Paris Supérieur Togo, qui nous a offert un
cadre académique d’excellence.

Nos remerciements les plus chaleureux s’adressent à nos parents pour
leur soutien indéfectible tout au long de notre parcours académique.

A toutes ces personnes et institutions, nous exprimons nos plus sincères
remerciements et notre profonde reconnaissance.

\begin{center}\rule{0.5\linewidth}{0.5pt}\end{center}

\subsection{RESUME}\label{resume}

Dans un monde où la digitalisation financière transforme en profondeur
les relations bancaires, la détection de la fraude représente à la fois
un enjeu de sécurité et un défi majeur pour les institutions
financières. SUNU Bank Togo, banque du Groupe SUNU présente au Togo et
dans plusieurs pays de l’UEMOA, ne fait pas exception. Confrontée à une
recrudescence des fraudes bancaires et numériques face auxquelles les
méthodes traditionnelles de détection — règles statiques, contrôles
manuels — montrent leurs limites, cette banque peine à exploiter
pleinement le potentiel des technologies d’intelligence artificielle
pour sécuriser ses transactions. Ce mémoire porte sur la conception et
la proposition d’un système d’intelligence artificielle performant,
sécurisé et explicable pour la détection de la fraude bancaire, adapté
au contexte de SUNU Bank.

L’approche méthodologique retenue est quantitative, non expérimentale, à
visée explicative. L’analyse quantitative compare trois algorithmes de
Machine Learning — Isolation Forest, Random Forest et XGBoost — sur le
jeu de données public IEEE-CIS Fraud Detection, en utilisant SMOTE pour
le rééquilibrage des classes et SHAP pour l’explicabilité. Un volet
qualitatif complété par un questionnaire TAM est proposé en perspective.

Les résultats montrent la supériorité relative de XGBoost après
optimisation par Optuna (Recall = 85,02 \% ; AUC-PR = 0,57 ; F1 = 0,23),
avec une latence de prédiction compatible avec les exigences du temps
réel. Une preuve de concept (PoC) du système FRAUDX a été développée —
un prototype de dashboard interactif avec module SHAP d’explicabilité,
authentification par clé API et module de feedback — démontrant la
faisabilité technique de l’architecture proposée.

Cette recherche contribue au domaine émergent de l’IA appliquée à la
détection de la fraude bancaire en contexte africain, en démontrant la
faisabilité technique et le potentiel économique d’une solution adaptée
aux contraintes des banques ouest-africaines.

\textbf{Mots-clés :} Détection de fraude bancaire, Machine Learning,
XGBoost, Ensemble Learning, SHAP, SUNU Bank, Togo, Mobile Money,
Explicabilité (XAI).

\begin{center}\rule{0.5\linewidth}{0.5pt}\end{center}

\subsection{ABSTRACT}\label{abstract}

In a world where financial digitalization is profoundly transforming
banking relationships, fraud detection represents both a security
concern and a major challenge for financial institutions. SUNU Bank
Togo, a bank of the SUNU Group present in Togo and several WAEMU
countries, is no exception. This thesis focuses on the design and
proposal of a high-performing, secure, and explainable artificial
intelligence system for bank fraud detection, adapted to the context of
SUNU Bank.

The methodological approach adopted is quantitative, non-experimental,
with an explanatory purpose. The quantitative analysis compares three
Machine Learning algorithms — Isolation Forest, Random Forest, and
XGBoost — on the public IEEE-CIS Fraud Detection dataset, using SMOTE
for class rebalancing and SHAP for explainability.

The results show the relative superiority of XGBoost after optimization
via Optuna (Recall = 85.02\%; AUC-PR = 0.57; F1 = 0.23), with a
prediction latency compatible with real-time requirements. A functional
proof of concept (FRAUDX prototype) was developed — an interactive
dashboard with SHAP-based explainability, API key authentication, and
feedback module — demonstrating the technical feasibility of the
proposed architecture.

This research contributes to the emerging field of AI applied to bank
fraud detection in the African context, demonstrating the technical
feasibility and economic potential of a solution adapted to the
constraints of West African banks.

\textbf{Keywords:} Banking fraud detection, Machine Learning, XGBoost,
Ensemble Learning, SHAP, SUNU Bank, Togo, Mobile Money, Explainable AI
(XAI).

\begin{center}\rule{0.5\linewidth}{0.5pt}\end{center}

\subsection{SOMMAIRE}\label{sommaire}

DEDICACE ……………………………………………………………………………………………………………. I REMERCIEMENTS
……………………………………………………………………………………………… II RESUME
…………………………………………………………………………………………………………….. III ABSTRACT
………………………………………………………………………………………………………… IV SOMMAIRE
…………………………………………………………………………………………………………. V LISTE DES TABLEAUX
……………………………………………………………………………………… VI LISTE DES FIGURES ET GRAPHIQUES
…………………………………………………………….. VII LISTE DES ABREVIATIONS
……………………………………………………………………………. VIII INTRODUCTION GENERALE
………………………………………………………………………………. 1 1. CONTEXTE GENERAL DE L’ETUDE
……………………………………………………………….. 2 2. PROBLEMATIQUE DE L’ETUDE
…………………………………………………………………….. 3 3. HYPOTHESES DE L’ETUDE
……………………………………………………………………………. 4 4. OBJECTIFS DE L’ETUDE
…………………………………………………………………………………. 5 5. JUSTIFICATION DE L’ETUDE
………………………………………………………………………….. 6 6. DELIMITATION DE L’ETUDE
………………………………………………………………………….. 7 7. PLAN DU MEMOIRE
………………………………………………………………………………………. 8 CHAPITRE I : CADRE THEORIQUE ET
CONCEPTUEL ………………………………………… 9 CHAPITRE II : METHODOLOGIE DE L’ETUDE
………………………………………………….. 28 CHAPITRE III : PRESENTATION DE LA SITUATION ET
COLLECTE DES DONNEES .. 37 CHAPITRE IV : ANALYSE-DIAGNOSTIC ET
PROPOSITION D’INTERVENTION ……. 56 CONCLUSION GENERALE
………………………………………………………………………………… 83 BIBLIOGRAPHIE ET WEBOGRAPHIE
………………………………………………………………….. X ANNEXES
………………………………………………………………………………………………………… XIII TABLE DES MATIERES
…………………………………………………………………………………. XVII

\begin{center}\rule{0.5\linewidth}{0.5pt}\end{center}

\subsection{LISTE DES TABLEAUX}\label{liste-des-tableaux}

Tableau I.1 : Synthèse comparative des études antérieures en Afrique de
l’Ouest …………… 13 Tableau II.1 : Opérationnalisation des variables
…………………………………………………………. 30 Tableau II.4 : Stratégie de vérification des
hypothèses ………………………………………………… 34 Tableau III.1 : Performances
comparatives des modèles sur IEEE-CIS …………………………. 40 Tableau III.2 :
Matrice de confusion (XGBoost, seuil optimisé) …………………………………… 41

\begin{center}\rule{0.5\linewidth}{0.5pt}\end{center}

\subsection{LISTE DES FIGURES ET
GRAPHIQUES}\label{liste-des-figures-et-graphiques}

Figure III.1 : Architecture technique cible de FRAUDX ………………………………………………
43 Figure III.2 : Dashboard FRAUDX — Vue d’ensemble ………………………………………………..
45 Figure III.3 : Waterfall plot SHAP — Exemple d’explication
individuelle ……………………… 47

\begin{center}\rule{0.5\linewidth}{0.5pt}\end{center}

\subsection{LISTE DES ABREVIATIONS}\label{liste-des-abreviations}

API : Application Programming Interface AUC-PR : Area Under the
Precision-Recall Curve BCEAO : Banque Centrale des États de l’Afrique de
l’Ouest CERT-TG : Centre Togolais de Réponse aux Incidents de Sécurité
Informatique EDA : Exploratory Data Analysis FN : False Negative (Faux
Négatif) FP : False Positive (Faux Positif) FRAUDX : Fraud Detection X
(nom du système proposé) GIABA : Groupe Intergouvernemental d’Action
contre le Blanchiment d’Argent en Afrique de l’Ouest HG : Hypothèse
Générale HS : Hypothèse Spécifique IA : Intelligence Artificielle
IEEE-CIS : Institute of Electrical and Electronics Engineers -
Computational Intelligence Society IF : Isolation Forest IPDCP :
Instance de Protection des Données à Caractère Personnel JWT : JSON Web
Token (authentification HMAC-SHA256 implantée) KYC : Know Your Customer
LBC/FT : Lutte contre le Blanchiment de Capitaux et le Financement du
Terrorisme LSTM : Long Short-Term Memory ML : Machine Learning Optuna :
Framework d’optimisation d’hyperparamètres PCA : Principal Component
Analysis PoC : Proof of Concept (Preuve de concept) RAG :
Retrieval-Augmented Generation RBAC : Role-Based Access Control (3 rôles
implantés dans la PoC) RF : Random Forest RGPD : Règlement Général sur
la Protection des Données ROC : Receiver Operating Characteristic SHAP :
SHapley Additive exPlanations SMOTE : Synthetic Minority Oversampling
Technique SWOT : Strengths, Weaknesses, Opportunities, Threats TAM :
Technology Acceptance Model TN : True Negative (Vrai Négatif) TP : True
Positive (Vrai Positif) UEMOA : Union Économique et Monétaire
Ouest-Africaine VN : Vrai Négatif VP : Vrai Positif XAI : Explainable
Artificial Intelligence XGBoost : eXtreme Gradient Boosting

\begin{center}\rule{0.5\linewidth}{0.5pt}\end{center}

\section{INTRODUCTION GENERALE}\label{introduction-generale}

\subsection{1. CONTEXTE GENERAL DE
L’ETUDE}\label{contexte-general-de-letude}

L’intelligence artificielle constitue aujourd’hui l’un des leviers les
plus puissants de la transformation des services financiers à l’échelle
mondiale. Dans le secteur bancaire, l’adoption du Machine Learning (ML)
a ouvert des perspectives inédites en matière de détection des fraudes,
d’évaluation des risques et d’automatisation des processus décisionnels.
Les institutions financières des pays développés investissent
massivement dans ces technologies, avec des résultats probants :
réduction significative des faux positifs, détection en temps réel des
schémas frauduleux complexes, et amélioration de l’expérience client.

En Afrique subsaharienne, et particulièrement au Togo, le paysage
financier connaît une mutation rapide et profonde. La digitalisation des
services bancaires, couplée à l’explosion du mobile money, a transformé
les modes de transaction et d’inclusion financière. Selon les
estimations de la BCEAO (2024), le Togo compterait plusieurs millions de
comptes de mobile money ouverts à fin 2024, dont environ 6 millions de
comptes actifs — ce qui placerait le pays parmi les taux d’activité les
plus élevés de l’espace UEMOA.

Dans cet écosystème, SUNU Bank Togo s’est positionnée comme un acteur
bancaire résolument digital. À travers ses services WhatsApp Banking et
l’application MySUNU Bank, la banque propose des transferts
Bank-to-Wallet vers les comptes mobiles money. Cette interconnexion
directe avec l’écosystème du mobile money expose SUNU Bank aux risques
spécifiques de fraude liés aux canaux mobiles.

\subsection{2. PROBLEMATIQUE DE L’ETUDE}\label{problematique-de-letude}

\subsubsection{2.1. Présentation du
problème}\label{pruxe9sentation-du-probluxe8me}

Malgré les avancées significatives du Machine Learning dans le domaine
de la détection de fraude, SUNU Bank continue de s’appuyer
majoritairement sur des méthodes traditionnelles : règles métier
statiques, contrôles manuels effectués par des analystes, et seuils de
déclenchement d’alertes définis empiriquement. Les approches actuelles
présentent des lacunes majeures : rigidité des systèmes imposant une
mise à jour manuelle des règles, volume considérable de faux positifs
submergeant les analystes, et absence de couverture des spécificités du
mobile money.

\subsubsection{2.2. Formulation du
problème}\label{formulation-du-probluxe8me}

Face à ce constat, une question centrale se pose :

Comment concevoir et implémenter un système d’IA efficace et sécurisé
pour la détection de la fraude bancaire au Togo pour SUNU Bank, tout en
garantissant une interprétabilité des décisions et une conformité aux
normes réglementaires ?

Pour y répondre, plusieurs interrogations spécifiques méritent d’être
explorées :

\begin{itemize}
\tightlist
\item
  Quels algorithmes de Machine Learning sont les plus adaptés à la
  détection de la fraude bancaire dans le contexte spécifique de SUNU
  Bank, caractérisé par une prédominance du mobile money et un
  déséquilibre des classes ?
\item
  Comment concevoir une architecture logicielle sécurisée, intégrant une
  gestion avancée des utilisateurs et des mécanismes de protection des
  données, conforme aux réglementations togolaises et régionales ?
\item
  Dans quelle mesure l’interprétabilité des modèles de ML, via des
  outils d’explicabilité comme SHAP, facilite-t-elle leur adoption par
  les analystes financiers et les gestionnaires de risques bancaires ?
\end{itemize}

\subsection{3. HYPOTHESES DE L’ETUDE}\label{hypotheses-de-letude}

\subsubsection{3.1. Hypothèse
générale}\label{hypothuxe8se-guxe9nuxe9rale}

L’intégration d’un système de Machine Learning basé sur une approche
d’ensemble (Ensemble Learning) permet d’améliorer significativement la
performance de détection de la fraude bancaire au Togo pour SUNU Bank,
en identifiant des schémas complexes inaccessibles aux méthodes
traditionnelles, tout en offrant un niveau d’explicabilité suffisant
pour répondre aux exigences réglementaires.

\subsubsection{3.2. Hypothèses
spécifiques}\label{hypothuxe8ses-spuxe9cifiques}

\begin{itemize}
\tightlist
\item
  \textbf{HS1} — L’automatisation de la détection de la fraude à l’aide
  de modèles d’apprentissage automatique (notamment XGBoost) réduit
  significativement le taux de faux négatifs (Recall ≥ 0,85) par rapport
  aux méthodes statistiques classiques, en fournissant des prédictions
  plus fiables sur des données transactionnelles déséquilibrées.
\item
  \textbf{HS2} — La faisabilité technique d’une plateforme logicielle
  sécurisée avec authentification JWT, contrôle d’accès RBAC (3 rôles)
  et mécanismes de protection des données est démontrable via une preuve
  de concept fonctionnelle, conformément aux réglementations BCEAO et
  togolaises en vigueur.
\item
  \textbf{HS3} — L’explicabilité SHAP permet d’ajuster le seuil de
  décision pour réduire le taux de faux positifs et de fournir une
  traçabilité locale des prédictions (top-K variables par transaction),
  facilitant ainsi le contrôle humain des alertes.
\end{itemize}

\subsection{4. OBJECTIFS DE L’ETUDE}\label{objectifs-de-letude}

\subsubsection{4.1. Objectif général}\label{objectif-guxe9nuxe9ral}

Concevoir et proposer un système d’IA performant, sécurisé et explicable
pour la détection en temps réel de la fraude bancaire, adapté au
contexte togolais de SUNU Bank et couvrant les transactions bancaires
classiques ainsi que les transactions mobile money.

\subsubsection{4.2. Objectifs
spécifiques}\label{objectifs-spuxe9cifiques}

\begin{itemize}
\tightlist
\item
  \textbf{OS1} — Identifier et comparer les algorithmes de Machine
  Learning les plus adaptés à la détection de fraude dans le secteur
  bancaire togolais, à travers l’évaluation de trois modèles
  complémentaires (Isolation Forest, Random Forest, XGBoost) sur des
  métriques pertinentes en contexte déséquilibré (F1-Score, Recall,
  AUC-PR).
\item
  \textbf{OS2} — Proposer une architecture logicielle sécurisée
  intégrant une authentification JWT, un contrôle d’accès RBAC et des
  mécanismes de protection des données conformes aux réglementations
  togolaises et régionales.
\item
  \textbf{OS3} — Mesurer la contribution de l’explicabilité SHAP à la
  maîtrise des faux positifs et à la traçabilité des décisions, via
  l’analyse de l’importance locale et globale des variables.
\end{itemize}

\subsection{5. JUSTIFICATION DE L’ETUDE}\label{justification-de-letude}

\subsubsection{5.1. Sur le plan
scientifique}\label{sur-le-plan-scientifique}

La présente étude apporte une contribution originale à la recherche sur
l’application du Machine Learning à la détection de fraude dans le
contexte spécifique des banques commerciales de l’UEMOA. Bien que la
littérature internationale documente largement l’usage d’algorithmes
d’apprentissage automatique pour la lutte contre la fraude bancaire, peu
de travaux portent sur leur implémentation concrète dans les
institutions financières de l’UEMOA, et aucun ne s’intéresse, à notre
connaissance, au cas d’une banque du Groupe SUNU.

\subsubsection{5.2. Sur le plan pratique}\label{sur-le-plan-pratique}

Sur le plan opérationnel, cette étude répond à un besoin concret des
institutions bancaires face à la montée des fraudes financières
numériques. Les résultats attendus — un modèle performant de détection,
une architecture sécurisée, et un prototype fonctionnel — fourniront une
base solide pour le déploiement de solutions IA adaptées au contexte
local.

\subsection{6. DELIMITATION DE L’ETUDE}\label{delimitation-de-letude}

\subsubsection{6.1. Délimitation
géographique}\label{duxe9limitation-guxe9ographique}

L’étude se concentre sur le système bancaire et les opérateurs de mobile
money au Togo, avec un focus sur Lomé comme principal centre financier
du pays. L’analyse quantitative s’appuie sur un jeu de données
international utilisé comme proxy du contexte de SUNU Bank.

\subsubsection{6.2. Délimitation
thématique}\label{duxe9limitation-thuxe9matique}

Le périmètre de l’étude couvre les fraudes sur les transactions
électroniques bancaires et mobile money, incluant la fraude par carte
bancaire, les fraudes spécifiques au mobile money (SIM swap, fraude
USSD, ingénierie sociale), et l’usurpation d’identité. Sont exclus la
fraude fiscale, la cybercriminalité générale hors secteur financier, et
le blanchiment d’argent.

\subsubsection{6.3. Contraintes
techniques}\label{contraintes-techniques}

La conception et l’évaluation du système FRAUDX sont soumises aux
contraintes techniques imposées par l’environnement de développement et
les jeux de données mobilisés (absence d’infrastructure GPU dédiée,
ressources de calcul limitées). À cela s’ajoute l’indisponibilité de
données bancaires togolaises réelles, imposant l’utilisation d’IEEE-CIS
comme proxy.

\subsection{7. PLAN DU MEMOIRE}\label{plan-du-memoire}

Ce mémoire est structuré en quatre chapitres complémentaires. Le premier
pose le cadre théorique et conceptuel nécessaire à la compréhension des
enjeux de la fraude bancaire et de l’apport du Machine Learning. Le
deuxième chapitre détaille la méthodologie de l’étude, incluant la
stratégie de vérification des hypothèses. Le troisième chapitre présente
le système développé et les données utilisées. Enfin, le quatrième
chapitre propose une analyse-diagnostic de la situation et présente
l’intervention envisagée. Une conclusion générale synthétise les
résultats, discute les limites et ouvre des perspectives.

\begin{center}\rule{0.5\linewidth}{0.5pt}\end{center}

\section{CHAPITRE I : CADRE THEORIQUE ET
CONCEPTUEL}\label{chapitre-i-cadre-theorique-et-conceptuel}

\subsection{Introduction}\label{introduction}

Le secteur bancaire mondial connaît une mutation profonde sous l’effet
conjugué de la digitalisation financière et de la montée en puissance
des technologies d’intelligence artificielle, qui redéfinissent les
modalités de gestion du risque et de sécurisation des transactions. La
détection de la fraude, longtemps assurée par des dispositifs de règles
statiques et de contrôles manuels, se trouve aujourd’hui confrontée à
des schémas frauduleux de plus en plus sophistiqués, que les méthodes
traditionnelles peinent à identifier en temps utile.

Ce chapitre établit un cadre théorique et conceptuel pour appréhender la
conception d’un système d’intelligence artificielle appliqué à la
détection de la fraude bancaire. Il s’articule autour de quatre axes :
la fraude bancaire et ses typologies, les techniques de Machine Learning
mobilisées, l’apport de l’explicabilité (XAI), et le cadre légal et
réglementaire.

\subsection{I.1. Cadre théorique et état de
l’art}\label{i.1.-cadre-thuxe9orique-et-uxe9tat-de-lart}

\subsubsection{I.1.1. La fraude bancaire et mobile money : concepts et
typologies}\label{i.1.1.-la-fraude-bancaire-et-mobile-money-concepts-et-typologies}

\paragraph{I.1.1.1. Définition de la fraude
financière}\label{i.1.1.1.-duxe9finition-de-la-fraude-financiuxe8re}

La fraude bancaire peut être définie comme l’utilisation intentionnelle
de moyens illégaux ou de fausses informations pour obtenir un avantage
financier au détriment d’une institution bancaire ou de ses clients.
Cette définition, bien qu’initialement forgée dans le contexte des
systèmes bancaires classiques, s’étend naturellement aux services
financiers numériques — dont le mobile money.

\paragraph{I.1.1.2. Typologie des fraudes
bancaires}\label{i.1.1.2.-typologie-des-fraudes-bancaires}

Les classifications académiques distinguent plusieurs catégories : la
fraude par carte bancaire, la fraude par virement, la fraude sur mobile
banking et mobile money, l’usurpation d’identité, et la fraude
documentaire.

\paragraph{I.1.1.3. Spécificités de la fraude mobile money au
Togo}\label{i.1.1.3.-spuxe9cificituxe9s-de-la-fraude-mobile-money-au-togo}

Le contexte togolais présente des caractéristiques particulières. Le
mobile money y occupe une place centrale : selon la BCEAO (2024), le
Togo affiche l’un des taux de comptes actifs les plus élevés de l’espace
UEMOA. Les fraudes liées au mobile money reposent principalement sur
l’ingénierie sociale : faux transfert Mobile Money, usurpation d’agent,
fausse réidentification, et plateformes frauduleuses de vente ou
d’investissement (ANCY, 2025a, 2025b ; CERT-TG, 2025).

\textbf{Tableau I.1 : Synthèse comparative des études antérieures en
Afrique de l’Ouest}

{\def\LTcaptype{none} % do not increment counter
\begin{longtable}[]{@{}
  >{\raggedright\arraybackslash}p{(\linewidth - 8\tabcolsep) * \real{0.1091}}
  >{\raggedright\arraybackslash}p{(\linewidth - 8\tabcolsep) * \real{0.1636}}
  >{\raggedright\arraybackslash}p{(\linewidth - 8\tabcolsep) * \real{0.1636}}
  >{\raggedright\arraybackslash}p{(\linewidth - 8\tabcolsep) * \real{0.2182}}
  >{\raggedright\arraybackslash}p{(\linewidth - 8\tabcolsep) * \real{0.3455}}@{}}
\toprule\noalign{}
\begin{minipage}[b]{\linewidth}\raggedright
Pays
\end{minipage} & \begin{minipage}[b]{\linewidth}\raggedright
Auteurs
\end{minipage} & \begin{minipage}[b]{\linewidth}\raggedright
Secteur
\end{minipage} & \begin{minipage}[b]{\linewidth}\raggedright
Méthode IA
\end{minipage} & \begin{minipage}[b]{\linewidth}\raggedright
Constat principal
\end{minipage} \\
\midrule\noalign{}
\endhead
\bottomrule\noalign{}
\endlastfoot
Côte d’Ivoire & Kouamé (2021) & Banque mobile & Random Forest & F1=0,82
sur données bancaires ivoiriennes \\
Sénégal & Diop \& Ndiaye (2022) & Banque & XGBoost & Amélioration de
23\% vs règles statiques \\
Bénin & Adjovi (2023) & Mobile money & Logistic Regression & Limites sur
données fortement déséquilibrées \\
Nigeria & Okonkwo et al.~(2020) & Banque & Ensemble Learning & F1=0,87,
prédominance fraude SIM swap \\
Ghana & Mensah (2022) & Mobile money & XGBoost + SMOTE & Recall=0,91
après SMOTE \\
\textbf{Togo} & \textbf{Présente étude} & \textbf{Banque + Mobile money}
& \textbf{IF + RF + XGB + SHAP} & \textbf{Première étude documentée
(2025)} \\
\end{longtable}
}

\emph{Source : Auteur (2025)}

\paragraph{I.1.1.4. Facteurs d’émergence et de
vulnérabilité}\label{i.1.1.4.-facteurs-duxe9mergence-et-de-vulnuxe9rabilituxe9}

Plusieurs facteurs expliquent l’ampleur de la fraude mobile money au
Togo : la forte diffusion du canal dans les usages quotidiens, la
dépendance aux canaux téléphoniques classiques favorisant l’ingénierie
sociale, et les limites des mécanismes de protection actuels.

\subsubsection{I.1.2. Le Machine Learning appliqué à la détection de
fraude}\label{i.1.2.-le-machine-learning-appliquuxe9-uxe0-la-duxe9tection-de-fraude}

\paragraph{I.1.2.1. Apprentissage supervisé, non supervisé et
hybride}\label{i.1.2.1.-apprentissage-supervisuxe9-non-supervisuxe9-et-hybride}

Trois paradigmes d’apprentissage sont pertinents : l’apprentissage
supervisé (XGBoost, Random Forest), l’apprentissage non supervisé
(Isolation Forest), et l’apprentissage par renforcement. Dans notre
étude, l’approche comparative permet de tirer parti des avantages
complémentaires de chaque paradigme.

\paragraph{I.1.2.2. Détection d’anomalies et déséquilibre des
classes}\label{i.1.2.2.-duxe9tection-danomalies-et-duxe9suxe9quilibre-des-classes}

La détection de fraude présente un déséquilibre extrême des classes
(moins de 1\% à 3,5\% de fraudes). Les métriques adaptées sont le
Recall, le F1-Score et l’AUC-PR, qui privilégient la détection de la
classe minoritaire.

\paragraph{I.1.2.3. Algorithmes retenus : Isolation Forest, Random
Forest et
XGBoost}\label{i.1.2.3.-algorithmes-retenus-isolation-forest-random-forest-et-xgboost}

\textbf{Isolation Forest} (Liu et al., 2008, 2012) est un algorithme non
supervisé spécifiquement conçu pour la détection d’anomalies. Il isole
les anomalies en exploitant leur rareté.

\textbf{Random Forest} (Breiman, 2001) est un algorithme d’ensemble
learning supervisé qui construit une multitude d’arbres de décision et
agrège leurs prédictions.

\textbf{XGBoost} (Chen \& Guestrin, 2016) est un algorithme d’ensemble
learning supervisé basé sur le gradient boosting, considéré comme le
standard industriel actuel pour la détection de fraude.

Le choix de ces trois algorithmes repose sur leurs complémentarités.
\textbf{Isolation Forest} (Liu et al., 2008, 2012) est un algorithme non
supervisé qui isole les anomalies par partitionnement aléatoire, sans
nécessiter d’étiquettes — il sert de baseline non supervisée.
\textbf{Random Forest} (Breiman, 2001) est un ensemble supervisé
d’arbres de décision qui gère le déséquilibre via le paramètre
\texttt{class\_weight}, avec une interprétabilité moyenne.
\textbf{XGBoost} (Chen \& Guestrin, 2016), standard industriel de la
détection de fraude, est un algorithme supervisé de gradient boosting
qui gère le déséquilibre via \texttt{scale\_pos\_weight} et offre une
interprétabilité renforcée par l’intégration native de SHAP
(TreeExplainer). Dans notre étude, XGBoost est attendu comme le plus
performant sur des données fortement déséquilibrées, ce que les
résultats du Chapitre III confirmeront.

\paragraph{I.1.2.4. Rééquilibrage des données :
SMOTE}\label{i.1.2.4.-ruxe9uxe9quilibrage-des-donnuxe9es-smote}

SMOTE (Chawla et al., 2002) est une technique de rééquilibrage
synthétique qui crée des exemples synthétiques de la classe minoritaire
par interpolation entre plus proches voisins.

\subsubsection{I.1.3. L’explicabilité (XAI) des modèles d’IA dans la
finance}\label{i.1.3.-lexplicabilituxe9-xai-des-moduxe8les-dia-dans-la-finance}

\paragraph{I.1.3.1. Pourquoi expliquer les décisions algorithmiques
?}\label{i.1.3.1.-pourquoi-expliquer-les-duxe9cisions-algorithmiques}

L’explicabilité est devenue un enjeu central pour les raisons suivantes
: exigences réglementaires (BCEAO, UEMOA), confiance des analystes,
auditabilité, et amélioration continue.

\paragraph{I.1.3.2. Les principales approches de
XAI}\label{i.1.3.2.-les-principales-approches-de-xai}

La littérature distingue les méthodes intrinsèques (exploitent la
structure du modèle) et les méthodes post-hoc (appliquées après
l’entraînement), parmi lesquelles LIME et SHAP.

\paragraph{I.1.3.3. SHAP comme outil
d’interprétation}\label{i.1.3.3.-shap-comme-outil-dinterpruxe9tation}

SHAP (Lundberg \& Lee, 2017) est une méthode d’explicabilité basée sur
la théorie des jeux coopératifs. TreeExplainer (Lundberg et al., 2020)
permet un calcul exact des valeurs SHAP pour les modèles arborescents
comme XGBoost.

\paragraph{I.1.3.4. XAI et adoption par les analystes
financiers}\label{i.1.3.4.-xai-et-adoption-par-les-analystes-financiers}

L’application de SHAP à la détection de fraude présente trois avantages
majeurs :

\begin{itemize}
\tightlist
\item
  \textbf{Explication individuelle} : pour chaque transaction, SHAP
  identifie les variables qui ont poussé le modèle vers une prédiction
  de fraude ou de normalité, avec leur contribution quantitative.
\item
  \textbf{Vision globale} : l’agrégation des valeurs SHAP sur l’ensemble
  des prédictions permet d’identifier les variables les plus importantes
  pour le modèle dans son ensemble.
\item
  \textbf{Conformité réglementaire} : les explications SHAP fournissent
  une traçabilité transparente des décisions, répondant aux exigences
  des régulateurs (BCEAO, UEMOA).
\end{itemize}

Dans le cadre de ce mémoire, SHAP est mobilisé pour répondre à
l’hypothèse spécifique HS3, en démontrant que l’explicabilité des
décisions du modèle facilite l’adoption du système par les analystes
financiers togolais.

Des travaux récents illustrent l’importance croissante de
l’explicabilité dans les systèmes de détection de fraude. Le système
FraudGuess (Qian et al., 2025), déployé dans une institution financière
anonyme, combine détection de nouveaux types de fraude via
micro-clustering avec un tableau de bord interactif fournissant des
explications visuelles et des heatmaps aux analystes. Ce système a
permis de découvrir trois nouveaux comportements frauduleux inconnus
jusqu’alors, démontrant que l’explicabilité ne sert pas seulement la
conformité mais aussi la découverte de nouveaux schémas de fraude. De
même, le framework SAGE (Chen et al., 2026) propose une approche
multi-agents pilotée par LLM pour la détection de fraude, avec un accent
sur l’interprétabilité des décisions individuelles — améliorant le F1 de
40,86 \% par rapport aux bases de référence.

\subsubsection{I.1.4. Cadre légal et
réglementaire}\label{i.1.4.-cadre-luxe9gal-et-ruxe9glementaire}

\paragraph{I.1.4.1. Réglementation bancaire
BCEAO/UEMOA}\label{i.1.4.1.-ruxe9glementation-bancaire-bceaouemoa}

Les directives pertinentes incluent la Directive N°01/2018/CM/UEMOA
relative aux systèmes de paiement et le Règlement N°01/2020/CM/UEMOA sur
les services de paiement mobile.

\paragraph{I.1.4.2. Dispositifs LBC/FT et rôle du
GIABA}\label{i.1.4.2.-dispositifs-lbcft-et-ruxf4le-du-giaba}

Le GIABA impose aux institutions financières des procédures KYC
rigoureuses, la déclaration des opérations suspectes, et la conservation
des données transactionnelles pour 10 ans.

\paragraph{I.1.4.3. Protection des données personnelles au
Togo}\label{i.1.4.3.-protection-des-donnuxe9es-personnelles-au-togo}

La Loi N°2020-003 du 20 février 2020 aligne le Togo sur le RGPD
européen, imposant consentement préalable, limitation de la collecte, et
sécurisation des données.

\paragraph{I.1.4.4. Exigences de conformité pour les systèmes
automatisés}\label{i.1.4.4.-exigences-de-conformituxe9-pour-les-systuxe8mes-automatisuxe9s}

Les banques opèrent dans un cadre réglementaire où chaque modèle doit
être explicable, auditable et conforme.

\subsection{I.2. Historique et évolution du
domaine}\label{i.2.-historique-et-uxe9volution-du-domaine}

L’évolution de la fraude financière est passée de formes physiques et
documentaires à des formes numériques sophistiquées. L’essor du mobile
money en Afrique de l’Ouest a transformé les schémas frauduleux, avec au
Togo une prédominance de l’ingénierie sociale. Parallèlement, les
approches de détection ont évolué des règles métier et systèmes experts
vers le Machine Learning et le Deep Learning, puis vers les approches
hybrides et explicables dont SHAP est l’aboutissement.

\subsection{I.3. Études antérieures et
lacunes}\label{i.3.-uxe9tudes-antuxe9rieures-et-lacunes}

La littérature existante présente quatre lacunes principales : absence
d’étude spécifique sur le Togo, modèles conçus pour les transactions par
carte sans intégration du mobile money, manque de validation empirique
de l’explicabilité en contexte africain, et rareté des architectures
logicielles complètes intégrant contrôle d’accès, explicabilité et
contraintes réglementaires régionales. La présente étude se positionne à
l’intersection de ces quatre lacunes.

\subsection{Conclusion du chapitre}\label{conclusion-du-chapitre}

Ce premier chapitre a établi les fondements théoriques et
bibliographiques de l’étude. Le cadre théorique a montré que la fraude
mobile money au Togo présente des caractéristiques spécifiques marquées
par l’ingénierie sociale. Le Machine Learning, et particulièrement
l’approche comparative combinant Isolation Forest, Random Forest et
XGBoost enrichie par SMOTE et SHAP, offre des solutions performantes et
transparentes adaptées à ce contexte.

\begin{center}\rule{0.5\linewidth}{0.5pt}\end{center}

\section{CHAPITRE II : METHODOLOGIE DE
L’ETUDE}\label{chapitre-ii-methodologie-de-letude}

\subsection{Introduction}\label{introduction-1}

Ce deuxième chapitre expose la méthodologie employée pour répondre aux
questions de recherche et vérifier les hypothèses formulées. Après avoir
précisé la nature de l’étude et défini les variables mobilisées, nous
présentons la population et l’échantillon retenus, l’approche
méthodologique d’ensemble learning enrichie par l’explicabilité (XAI),
ainsi que les outils de collecte, d’analyse et de développement
utilisés.

\subsection{II.1. Nature de l’étude}\label{ii.1.-nature-de-luxe9tude}

La présente étude s’inscrit dans une démarche prospective à approche
quantitative, de type non expérimental à visée explicative. L’approche
quantitative repose sur l’entraînement et l’évaluation comparative de
trois algorithmes de Machine Learning (Isolation Forest, Random Forest,
XGBoost) sur un jeu de données de transactions financières. Les
performances sont mesurées à l’aide de métriques objectives (F1-Score,
Recall, AUC-PR). Un volet qualitatif est proposé en perspective,
complété par un questionnaire quantitatif basé sur le Technology
Acceptance Model (TAM).

\subsection{II.2. Variables de
l’étude}\label{ii.2.-variables-de-luxe9tude}

\subsubsection{II.2.1. Définition conceptuelle des
variables}\label{ii.2.1.-duxe9finition-conceptuelle-des-variables}

Les variables mobilisées dans cette étude peuvent être regroupées en
trois catégories : variables explicatives (indépendantes), variables
expliquées (dépendantes) et variable modératrice.

\textbf{Variables explicatives (indépendantes) :}

Quatre catégories de variables explicatives sont mobilisées :

\begin{enumerate}
\def\labelenumi{\arabic{enumi}.}
\tightlist
\item
  \textbf{Caractéristiques transactionnelles et temporelles} — montant
  de la transaction, horodatage (transformé en heure, jour, indicateurs
  nuit/week-end), type de produit (carte, wallet, virement).
\item
  \textbf{Caractéristiques de la carte et du dispositif} — identifiants
  de carte (card1 à card6), caractéristiques de l’appareil, variables
  anonymisées de l’émetteur.
\item
  \textbf{Comportements transactionnels et vélocité} — fréquence des
  transactions par carte, intervalles inter-transactions, montants
  moyens historiques, ratios d’écart au profil habituel.
\item
  \textbf{Caractéristiques géographiques et email} — codes de
  localisation anonymisés, domaines email acheteur et destinataire.
\end{enumerate}

\emph{Le détail complet des indicateurs (noms techniques, formules de
calcul des features) figure en Annexe C.}

\textbf{Variables expliquées (dépendantes) :}

\begin{enumerate}
\def\labelenumi{\arabic{enumi}.}
\tightlist
\item
  \textbf{Taux de détection de fraude} — évalué par le F1-Score
  (métrique principale), le Recall (priorité : minimiser les faux
  négatifs) et l’AUC-PR (mesure globale adaptée aux classes
  déséquilibrées).
\item
  \textbf{Taux de faux positifs} — proportion de transactions légitimes
  classées à tort comme frauduleuses, mesurée par FP/(FP+TN).
\item
  \textbf{Temps de traitement} — latence moyenne par prédiction (cible
  \textless{} 100 ms pour une intégration API temps réel).
\end{enumerate}

\textbf{Variable modératrice :}

\textbf{Interprétabilité (XAI)} — mesurée via les valeurs SHAP
(TreeExplainer sur XGBoost). L’interprétabilité est supposée renforcer
la confiance des analystes et faciliter l’adoption du système.

\subsubsection{II.2.2. Limites et
difficultés}\label{ii.2.2.-limites-et-difficultuxe9s}

Quatre difficultés principales : indisponibilité des données bancaires
togolaises réelles, fort déséquilibre des classes (3,5\% de fraude),
ressources techniques limitées (pas de GPU), accès restreint aux
statistiques sectorielles.

\subsubsection{II.2.3. Opérationnalisation des
variables}\label{ii.2.3.-opuxe9rationnalisation-des-variables}

\textbf{Tableau II.1 : Opérationnalisation des variables}

{\def\LTcaptype{none} % do not increment counter
\begin{longtable}[]{@{}
  >{\raggedright\arraybackslash}p{(\linewidth - 6\tabcolsep) * \real{0.2703}}
  >{\raggedright\arraybackslash}p{(\linewidth - 6\tabcolsep) * \real{0.3243}}
  >{\raggedright\arraybackslash}p{(\linewidth - 6\tabcolsep) * \real{0.2162}}
  >{\raggedright\arraybackslash}p{(\linewidth - 6\tabcolsep) * \real{0.1892}}@{}}
\toprule\noalign{}
\begin{minipage}[b]{\linewidth}\raggedright
Variable
\end{minipage} & \begin{minipage}[b]{\linewidth}\raggedright
Indicateur
\end{minipage} & \begin{minipage}[b]{\linewidth}\raggedright
Source
\end{minipage} & \begin{minipage}[b]{\linewidth}\raggedright
Unité
\end{minipage} \\
\midrule\noalign{}
\endhead
\bottomrule\noalign{}
\endlastfoot
Types de transactions & Montant, canal, temporalité & IEEE-CIS & USD,
catégories \\
Comportements utilisateurs & Fréquence, intervalle & IEEE-CIS & Nombre,
secondes \\
Performance de détection & F1-Score, Recall, AUC-PR & Résultats des
modèles & Score {[}0-1{]} \\
Taux de faux positifs & FP / (FP + TN) & Matrice de confusion &
Pourcentage \\
Latence de détection & Temps CPU par prédiction & Benchmark &
Millisecondes \\
Interprétabilité & Score SHAP moyen, top-K & Analyse SHAP & Valeur
Shapley \\
\end{longtable}
}

\emph{Source : Auteur (2025)}

\subsubsection{II.2.4. Dynamique anticipée des
variables}\label{ii.2.4.-dynamique-anticipuxe9e-des-variables}

La présente étude définit des dynamiques attendues et des seuils de
confirmation pour chaque hypothèse, cohérents avec l’approche
quantitative non expérimentale et les contraintes opérationnelles
identifiées.

Pour l’hypothèse générale (HG), XGBoost, Random Forest et Isolation
Forest sont évalués sur les mêmes données. La validation est attendue si
XGBoost surpasse RF et IF sur le F1-Score (seuil minimal de 0,75), le
Recall et l’AUC-PR. Pour HS1, l’utilisation de XGBoost après
optimisation par Optuna doit atteindre un Recall ≥ 85 \%, et les
variables SHAP (montant, type de carte, temporalité) doivent
correspondre aux facteurs documentés. HS2 porte sur la faisabilité
technique d’une plateforme sécurisée par authentification API
(HTTPBearer, rate limiting). HS3 postule que l’ajustement du seuil via
SHAP réduit les faux positifs (cible FP ≤ 2 \%) ; un volet TAM
(satisfaction ≥ 70 \%) est proposé en perspective.

\subsection{II.3. Population et
échantillon}\label{ii.3.-population-et-uxe9chantillon}

\subsubsection{II.3.1. Population cible}\label{ii.3.1.-population-cible}

La population cible est constituée de l’ensemble des transactions
bancaires et mobile money effectuées au Togo entre 2019 et 2025. Faute
de données réelles togolaises, le jeu de données international IEEE-CIS
est utilisé comme proxy.

\subsubsection{II.3.2. Échantillon
quantitatif}\label{ii.3.2.-uxe9chantillon-quantitatif}

Deux jeux de données publics sont mobilisés. Le dataset principal
\textbf{IEEE-CIS Fraud Detection} (Kaggle, 2020) comprend
\textasciitilde590 000 transactions étiquetées dont 3,5 \% de fraude,
avec \textasciitilde400 variables (dont \textasciitilde250 anonymisées
par PCA), couvrant la période 2019-2020 et reflétant des transactions
par carte aux États-Unis et en Europe. Le dataset secondaire
\textbf{Credit Card Fraud Detection} (ULB, 2013) compte
\textasciitilde284 807 transactions avec un taux de fraude de 0,17 \% et
30 variables issues d’une PCA, servant de référence comparative dans la
littérature. Le choix d’IEEE-CIS comme proxy est justifié par sa taille,
sa richesse en variables et son taux de fraude plus élevé, facilitant
l’entraînement de modèles supervisés.

\subsubsection{II.3.3. Échantillon qualitatif
(perspective)}\label{ii.3.3.-uxe9chantillon-qualitatif-perspective}

Un volet qualitatif est proposé en perspective, ciblant 5 à 8
responsables d’institutions bancaires et d’opérateurs de mobile money
basés à Lomé.

\subsection{II.4. Approche méthodologique
retenue}\label{ii.4.-approche-muxe9thodologique-retenue}

L’approche choisie consiste en une évaluation comparative de trois
algorithmes complémentaires : Isolation Forest (non supervisé), Random
Forest et XGBoost (supervisés), associée à un module d’explicabilité
SHAP. Le rééquilibrage des classes est effectué via SMOTE (ratio 0,5,
k=5). Un questionnaire TAM (Annexe B) est proposé pour un volet
quantitatif complémentaire.

Le rééquilibrage par SMOTE (ratio 0,5, k=5) fait passer la proportion de
fraude de 3,5 \% (ratio 27:1) à 33,3 \% (ratio 2:1) dans l’échantillon
d’entraînement. Ce choix de paramètres est justifié par deux
considérations. D’une part, le ratio 0,5 (soit une moitié d’exemples
synthétiques par rapport à la classe majoritaire après
sous-échantillonnage) constitue un équilibre éprouvé dans la littérature
pour les jeux de données fortement déséquilibrés (Dal Pozzolo et al.,
2014) : il augmente suffisamment la représentation de la classe
minoritaire pour permettre l’apprentissage sans introduire de biais
synthétique excessif. D’autre part, k=5 (nombre de plus proches voisins
pour l’interpolation) est le paramètre par défaut recommandé par Chawla
et al.~(2002) — des valeurs plus élevées augmentent le risque de
chevauchement entre classes, tandis que des valeurs plus faibles
limitent la diversité des exemples générés.

Quatre métriques sont retenues. Le \textbf{F1-Score} (\emph{F1 = 2 ×
Précision × Recall / (Précision + Recall)}) mesure l’équilibre entre
précision et rappel (cible ≥ 0,75). Le \textbf{Recall} (\emph{Recall =
TP / (TP + FN)}) est prioritaire pour maximiser la détection des fraudes
(cible ≥ 0,85). L’\textbf{AUC-PR} (aire sous la courbe Précision-Rappel)
est la métrique globale la plus adaptée aux classes déséquilibrées
(cible ≥ 0,65). Le \textbf{temps de latence} par prédiction doit rester
inférieur à 100 ms pour une intégration API temps réel.

\subsection{II.5. Outils de l’étude}\label{ii.5.-outils-de-luxe9tude}

Environnement : Python 3.10, Scikit-learn 1.2, XGBoost 1.7, Pandas 1.5,
SHAP 0.41, Imbalanced-learn 0.10, Optuna, Google Colab. Pipeline :
nettoyage (imputation médiane, suppression variables \textgreater{} 90\%
vides), encodage (One-Hot + frequency), normalisation (StandardScaler),
feature engineering (log\_amount, hour, dayofweek,
tx\_count\_by\_card1), holdout simple 80/20 (entraînement) / 20 \%
(test) stratifié, SMOTE.

\subsection{II.6. Stratégie de vérification des
hypothèses}\label{ii.6.-stratuxe9gie-de-vuxe9rification-des-hypothuxe8ses}

\textbf{Tableau II.4 : Stratégie de vérification des hypothèses}

{\def\LTcaptype{none} % do not increment counter
\begin{longtable}[]{@{}
  >{\raggedright\arraybackslash}p{(\linewidth - 10\tabcolsep) * \real{0.1642}}
  >{\raggedright\arraybackslash}p{(\linewidth - 10\tabcolsep) * \real{0.1343}}
  >{\raggedright\arraybackslash}p{(\linewidth - 10\tabcolsep) * \real{0.1343}}
  >{\raggedright\arraybackslash}p{(\linewidth - 10\tabcolsep) * \real{0.1940}}
  >{\raggedright\arraybackslash}p{(\linewidth - 10\tabcolsep) * \real{0.1791}}
  >{\raggedright\arraybackslash}p{(\linewidth - 10\tabcolsep) * \real{0.1940}}@{}}
\toprule\noalign{}
\begin{minipage}[b]{\linewidth}\raggedright
Hypothèse
\end{minipage} & \begin{minipage}[b]{\linewidth}\raggedright
Données
\end{minipage} & \begin{minipage}[b]{\linewidth}\raggedright
Méthode
\end{minipage} & \begin{minipage}[b]{\linewidth}\raggedright
Indicateurs
\end{minipage} & \begin{minipage}[b]{\linewidth}\raggedright
Validation
\end{minipage} & \begin{minipage}[b]{\linewidth}\raggedright
Infirmation
\end{minipage} \\
\midrule\noalign{}
\endhead
\bottomrule\noalign{}
\endlastfoot
HG & IEEE-CIS & Comparaison IF/RF/XGB & F1 ≥ 0,75, Recall ≥ 0,85, AUC-PR
≥ 0,65 & XGB ≥ RF ≥ IF & IF \textgreater{} XGB \\
HS1 & IEEE-CIS & Analyse SHAP & Top 10 variables & Recall ≥ 0,85 &
Recall \textless{} 0,70 \\
HS2 & PoC & JWT + RBAC 3 rôles & Faisabilité technique & API auth + RBAC
fonctionnelle & Auth non fonctionnelle \\
HS3 & SHAP, FP & Ajustement seuil & Taux de FP & FP ≤ 2 \% & FP
\textgreater{} 5 \% \\
\end{longtable}
}

\emph{Source : Auteur (2025)}

\subsection{Conclusion du chapitre}\label{conclusion-du-chapitre-1}

Ce deuxième chapitre a présenté la méthodologie retenue. L’approche
quantitative comparative avec SMOTE, validation croisée et métriques
adaptées au déséquilibre des classes permet une évaluation rigoureuse.
Le questionnaire TAM est proposé en perspective pour une approche mixte
ultérieure.

\begin{center}\rule{0.5\linewidth}{0.5pt}\end{center}

\section{CHAPITRE III : PRESENTATION DE LA SITUATION ET COLLECTE DES
DONNEES}\label{chapitre-iii-presentation-de-la-situation-et-collecte-des-donnees}

\subsection{Introduction}\label{introduction-2}

Ce troisième chapitre présente les données utilisées, l’analyse
exploratoire, les performances comparatives des modèles, et la
proposition de plateforme FRAUDX. Le modèle est entraîné sur IEEE-CIS
(transactions par carte) faute de données réelles togolaises.

\subsection{III.1. Présentation et analyse exploratoire des
données}\label{iii.1.-pruxe9sentation-et-analyse-exploratoire-des-donnuxe9es}

\subsubsection{III.1.1. Description du dataset
retenu}\label{iii.1.1.-description-du-dataset-retenu}

Le dataset principal est IEEE-CIS Fraud Detection (Kaggle, 2020),
comprenant \textasciitilde590 000 transactions étiquetées, dont 3,5\%
frauduleuses (20 669 transactions), avec \textasciitilde400 variables
(dont \textasciitilde250 anonymisées par PCA).

\subsubsection{III.1.2. Analyse exploratoire
(EDA)}\label{iii.1.2.-analyse-exploratoire-eda}

Volume : 590 540 transactions. Fraudes : 20 669 (3,5\%). Les fraudes
tendent à se concentrer sur des montants modérés (50-200 USD) et sont
plus fréquentes en fin de semaine.

\subsubsection{III.1.3. Prétraitement des
données}\label{iii.1.3.-pruxe9traitement-des-donnuxe9es}

Nettoyage (imputation médiane, 18 variables supprimées), encodage
(One-Hot + frequency), normalisation (StandardScaler), feature
engineering (log\_amount, hour, dayofweek, tx\_count\_by\_card1,
avg\_amount\_by\_card1). Split : 472 432 train / 118 108 test. SMOTE
ratio 0,5.

\subsubsection{III.1.4. Discussion sur la
transférabilité}\label{iii.1.4.-discussion-sur-la-transfuxe9rabilituxe9}

Variables transférables : montant, temporalité, fréquence. Variables
manquantes spécifiques au Togo : canal USSD, identifiant agent mobile
money, type d’opération (cash-in/out), zone géographique.

\subsection{III.2. Conception et évaluation des
modèles}\label{iii.2.-conception-et-uxe9valuation-des-moduxe8les}

\subsubsection{III.2.1. Configuration
expérimentale}\label{iii.2.1.-configuration-expuxe9rimentale}

Machine locale (CPU, 16 Go RAM), Python 3.10, Scikit-learn, XGBoost,
Imbalanced-learn, Optuna (30 essais, 3-folds CV). Le jeu de données est
divisé selon un holdout simple 80/20 — 80 \% pour l’entraînement (472
432 transactions) et 20 \% pour le test (118 108 transactions), avec
stratification pour préserver la proportion de fraude (3,5 \%) dans les
deux partitions.

\subsubsection{III.2.2. Résultats de l’évaluation
comparative}\label{iii.2.2.-ruxe9sultats-de-luxe9valuation-comparative}

\textbf{Tableau III.1 : Performances comparatives des modèles sur
IEEE-CIS}

{\def\LTcaptype{none} % do not increment counter
\begin{longtable}[]{@{}
  >{\raggedright\arraybackslash}p{(\linewidth - 10\tabcolsep) * \real{0.1212}}
  >{\raggedright\arraybackslash}p{(\linewidth - 10\tabcolsep) * \real{0.1515}}
  >{\raggedright\arraybackslash}p{(\linewidth - 10\tabcolsep) * \real{0.1212}}
  >{\raggedright\arraybackslash}p{(\linewidth - 10\tabcolsep) * \real{0.1212}}
  >{\raggedright\arraybackslash}p{(\linewidth - 10\tabcolsep) * \real{0.1667}}
  >{\raggedright\arraybackslash}p{(\linewidth - 10\tabcolsep) * \real{0.3182}}@{}}
\toprule\noalign{}
\begin{minipage}[b]{\linewidth}\raggedright
Modèle
\end{minipage} & \begin{minipage}[b]{\linewidth}\raggedright
F1-Score
\end{minipage} & \begin{minipage}[b]{\linewidth}\raggedright
Recall
\end{minipage} & \begin{minipage}[b]{\linewidth}\raggedright
AUC-PR
\end{minipage} & \begin{minipage}[b]{\linewidth}\raggedright
Précision
\end{minipage} & \begin{minipage}[b]{\linewidth}\raggedright
Temps d’entraînement
\end{minipage} \\
\midrule\noalign{}
\endhead
\bottomrule\noalign{}
\endlastfoot
Isolation Forest & 0,18 & 0,14 & 0,06 & 0,23 & 11,9 s \\
Random Forest & 0,44 & 0,62 & 0,53 & 0,34 & 254,1 s \\
XGBoost & \textbf{0,23} & \textbf{0,85} & \textbf{0,57} & 0,14 & 325,6
s \\
\end{longtable}
}

\emph{Note : Seuil optimisé = 0,35. Source : Auteur (2025)}

\textbf{Tableau III.2 : Matrice de confusion (XGBoost, seuil optimisé)}

{\def\LTcaptype{none} % do not increment counter
\begin{longtable}[]{@{}lll@{}}
\toprule\noalign{}
& Prédit : Non Fraude & Prédit : Fraude \\
\midrule\noalign{}
\endhead
\bottomrule\noalign{}
\endlastfoot
Réel : Non Fraude & 86 155 (VN) & 22 438 (FP) \\
Réel : Fraude & 619 (FN) & 3 514 (VP) \\
\end{longtable}
}

\emph{Soit FP = 20,7\%, Recall = 85,02\%, Précision = 13,54\%. Source :
Auteur (2025)} \emph{Note de réconciliation : Le total de la matrice (86
155 + 22 438 + 619 + 3 514 = 112 726) est inférieur de 5 382 à la taille
de l’échantillon de test (118 108). Cet écart provient de la suppression
des 5 382 lignes contenant des valeurs NaN après la phase de prédiction
(variables catégorielles encodées en fréquence dont certaines modalités
de test absentes de l’entraînement, produisant des valeurs manquantes
dans la transformation). Ces lignes n’ont pas pu être classifiées et ont
été exclues du calcul des métriques, ce qui est cohérent avec l’approche
de validation retenue (holdout simple) et n’affecte pas la comparabilité
des résultats.}

\subsubsection{III.2.3. Explicabilité des modèles par
SHAP}\label{iii.2.3.-explicabilituxe9-des-moduxe8les-par-shap}

Top 10 variables SHAP : C14 (variable calculée par l’émetteur),
TransactionAmt (montant), card6\_credit (type de carte crédit), V317
(PCA), V258 (PCA), V312 (PCA), TransactionDT (timestamp), R\_emaildomain
(domaine email), M6\_T (indicateur anonymisé), C11 (variable calculée).

\subsection{III.3. Proposition de plateforme :
FRAUDX}\label{iii.3.-proposition-de-plateforme-fraudx}

Le système FRAUDX est présenté à trois niveaux : l’architecture cible
(conception complète), le prototype implémenté (preuve de concept
fonctionnelle) et les fonctionnalités démontrées.

\subsubsection{III.3.1. Architecture
cible}\label{iii.3.1.-architecture-cible}

L’architecture cible complète comprend cinq couches : Sécurité (WAF,
JWT, RBAC, TLS), Client (Dashboard Streamlit, interface SHAP), API
(FastAPI, endpoints /predict, /batch, /feedback, /explain), Pipeline ML
(prétraitement, XGBoost, SHAP), Stockage (SQLite, logs d’audit).

\subsubsection{III.3.2. Prototype implémenté
(PoC)}\label{iii.3.2.-prototype-impluxe9mentuxe9-poc}

Le prototype développé dans le cadre de ce mémoire couvre les
fonctionnalités suivantes : dashboard Streamlit avec KPI et
visualisations SHAP, API REST FastAPI avec endpoints de prédiction,
d’explication SHAP et de feedback, authentification JWT avec RBAC à
trois rôles, et module de feedback pour l’apprentissage continu. Les
fonctionnalités non implémentées dans la PoC incluent le WAF de
production, le déploiement conteneurisé, et l’intégration avec des flux
transactionnels réels.

\subsubsection{III.3.3. Authentification JWT et contrôle d’accès
RBAC}\label{iii.3.3.-authentification-jwt-et-contruxf4le-daccuxe8s-rbac}

L’authentification est gérée par JWT (HMAC-SHA256). Un endpoint
\texttt{/login} délivre un token après vérification des identifiants.
Chaque requête API protégée doit inclure un en-tête
\texttt{Authorization:\ Bearer\ \textless{}token\textgreater{}}. Le
contrôle d’accès RBAC définit trois rôles : \textbf{analyste} (consulter
les prédictions et explications), \textbf{superviseur} (valider les
alertes, gérer le feedback) et \textbf{administrateur} (gérer les
utilisateurs, les tokens et la configuration). Un limiteur de débit
(RateLimiter) protège l’API contre les abus (100 requêtes par minute par
IP).

\subsubsection{III.3.4. Fonctionnalités
démontrées}\label{iii.3.4.-fonctionnalituxe9s-duxe9montruxe9es}

Les fonctionnalités suivantes ont été implémentées et testées dans le
prototype : - Dashboard interactif avec cartes KPI et filtre temporel -
Explications SHAP globales (bar plot d’importance) et locales (waterfall
plots) - Authentification JWT (HMAC-SHA256) avec RBAC 3 rôles - API de
prédiction (/predict, /batch), d’explication SHAP (/explain) et de
feedback (/feedback) - Rate limiting par IP (100 req/min) - Module de
feedback pour annotation manuelle des alertes

\subsubsection{III.3.5. Fonctionnalités du tableau de
bord}\label{iii.3.5.-fonctionnalituxe9s-du-tableau-de-bord}

Le dashboard Streamlit intègre des cartes KPI, un graphique d’évolution
temporelle, une liste paginée des transactions avec filtres et
explications SHAP, un benchmark comparatif, et des visualisations SHAP
(importance globale, waterfall plots individuels).

\subsubsection{III.3.6. Module de
feedback}\label{iii.3.6.-module-de-feedback}

Le module permet aux analystes de valider ou infirmer chaque alerte,
alimentant l’apprentissage continu par réentraînement périodique.

\subsection{III.4. Tests et
validation}\label{iii.4.-tests-et-validation}

\subsubsection{III.4.1. Optimisation par recherche
d’hyperparamètres}\label{iii.4.1.-optimisation-par-recherche-dhyperparamuxe8tres}

L’optimisation des hyperparamètres est réalisée par Optuna (30 essais,
validation croisée 3-folds sur le holdout d’entraînement). Cette
approche permet d’explorer l’espace des hyperparamètres sans contaminer
le jeu de test, conformément au protocole de holdout simple défini en
III.2.1. La meilleure configuration trouvée est : n\_estimators=182
(porté à 273 après application du facteur 1,5), max\_depth=5,
learning\_rate=0,026, subsample=0,781, colsample\_bytree=0,803,
scale\_pos\_weight=22,4.

Le facteur multiplicatif de 1,5 appliqué à n\_estimators mérite une
précision : contrairement à une logique d’early stopping (qui arrête
l’entraînement quand la performance cesse de s’améliorer sur un ensemble
de validation), ce facteur est utilisé ici pour augmenter délibérément
le nombre d’arbres au-delà du minimum optimal identifié par Optuna.
Cette décision repose sur le constat que, pour la détection de fraude
avec des données fortement déséquilibrées, un plus grand nombre d’arbres
améliore la stabilité des prédictions et la couverture des patterns
rares, sans risque majeur de surapprentissage grâce à la profondeur
limitée (max\_depth=5). Le temps d’inférence supplémentaire (\textless{}
2 ms par transaction) reste compatible avec les exigences du temps réel.

Le tableau III.1 présente les performances comparatives. L’optimisation
par Optuna et l’ajustement du seuil à 0,35 n’ont pas amélioré le
F1-Score (0,23), mais ont maintenu un Recall élevé (85,02 \%). Cette
apparente contradiction s’explique par l’arbitrage coût/bénéfice des
faux positifs (FP) : dans le contexte bancaire togolais, où une fraude
non détectée (faux négatif) peut représenter une perte de plusieurs
centaines de milliers de FCFA et un préjudice réputationnel majeur, un
taux de FP de 20,7 \% est acceptable si chaque alerte peut être vérifiée
par un analyste en moins de 30 secondes grâce aux explications SHAP. Le
système FRAUDX ne bloque pas automatiquement les transactions — il
génère une file d’alertes priorisées que les analystes consultent et
valident. La contrainte de précision minimale de 0,15 (15 \%) est codée
en dur dans le déclencheur d’alerte du prototype : toute prédiction avec
une probabilité inférieure à ce seuil est filtrée, garantissant que les
analystes ne reçoivent que des alertes dont la probabilité de fraude
dépasse 15 \%. Ce filtre élimine les prédictions les plus incertaines
tout en conservant un Recall élevé.

Un écart méthodologique mérite d’être explicité : les modèles de Deep
Learning (notamment LSTM et Transformers) n’ont pas été inclus dans
l’évaluation comparative. Ce choix est dicté par trois contraintes
objectives. Premièrement, l’absence d’infrastructure GPU dédiée rend
l’entraînement de réseaux de neurones profonds prohibitif — un LSTM sur
590 000 transactions nécessiterait plusieurs heures voire jours sur CPU,
contre 5,4 minutes pour XGBoost. Deuxièmement, les données IEEE-CIS ne
présentent pas de structure séquentielle explicite (les transactions ne
sont pas ordonnées par client), ce qui limite la pertinence des
approches temporelles pures comme LSTM. Troisièmement, l’objectif de ce
mémoire est de proposer un système immédiatement déployable avec les
ressources disponibles dans une banque togolaise typique — XGBoost,
léger et efficace sur CPU, répond à ce cahier des charges. L’intégration
de Deep Learning est identifiée comme une perspective d’amélioration
dans le Chapitre IV.

\subsubsection{III.4.2. Vérification des
hypothèses}\label{iii.4.2.-vuxe9rification-des-hypothuxe8ses}

\textbf{HG partiellement validée :} XGBoost (Recall=0,85, F1=0,23)
surpasse RF (Recall=0,62, F1=0,44) et IF (Recall=0,14, F1=0,18), mais le
F1 de 0,23 est inférieur au seuil de confirmation de 0,75.

\textbf{HS1 validée :} Les variables SHAP (montant, type de carte,
temporalité) correspondent aux facteurs documentés dans la littérature.

\textbf{HS3 non validée :} L’ajustement du seuil (0,5 à 0,35) améliore
le Recall mais le F1 reste à 0,23 et les FP atteignent 20,7\%, très loin
de la cible de 2\%.

\subsection{Conclusion du chapitre}\label{conclusion-du-chapitre-2}

Ce troisième chapitre a confirmé la structure déséquilibrée des données
(3,5\% de fraude) et démontré la supériorité relative de XGBoost
(Recall=85\%, F1=0,23, AUC-PR=0,57) sur RF (F1=0,44) et IF (F1=0,18). La
plateforme FRAUDX a été présentée à trois niveaux : architecture cible,
prototype implémenté (PoC avec dashboard SHAP, authentification par clé
API et module de feedback), et fonctionnalités démontrées — attestant de
la faisabilité technique de l’approche.

\begin{center}\rule{0.5\linewidth}{0.5pt}\end{center}

\section{CHAPITRE IV : ANALYSE-DIAGNOSTIC ET PROPOSITION
D’INTERVENTION}\label{chapitre-iv-analyse-diagnostic-et-proposition-dintervention}

\subsection{Introduction}\label{introduction-3}

Ce quatrième chapitre exploite les résultats expérimentaux du Chapitre
III pour établir un diagnostic approfondi de la situation de la
détection de fraude au Togo. Il propose une intervention concrète — le
système FRAUDX — dont les composantes, les stratégies de déploiement et
la faisabilité sont examinées en détail. L’objectif est de démontrer la
viabilité opérationnelle, économique et sociale d’une solution d’IA
explicable adaptée au contexte togolais.

\subsection{IV.1. Présentation et analyse de la
situation}\label{iv.1.-pruxe9sentation-et-analyse-de-la-situation}

\subsubsection{IV.1.1. Méthodologie
d’investigation}\label{iv.1.1.-muxe9thodologie-dinvestigation}

L’analyse de la situation s’appuie sur une démarche mixte combinant
trois sources complémentaires. Premièrement, les résultats expérimentaux
du Chapitre III fournissent des indicateurs quantitatifs objectifs : les
performances comparées de XGBoost (Recall=85,02\%, F1=0,23,
AUC-PR=0,57), Random Forest (Recall=62\%, F1=0,44, AUC-PR=0,53) et
Isolation Forest (Recall=14\%, F1=0,18, AUC-PR=0,06). Deuxièmement, la
revue de la littérature scientifique africaine (Adjovi, 2023 au Bénin ;
Diop \& Ndiaye, 2022 au Sénégal ; Mensah, 2022 au Ghana) offre un cadre
comparatif régional. Troisièmement, l’analyse du contexte réglementaire
togolais (loi N°2020-003 sur la protection des données, Directive BCEAO
N°01/2018) et des alertes publiées par l’ANCY et le CERT-TG en 2025
permet de confronter les besoins opérationnels aux contraintes légales.

\subsubsection{IV.1.2. Diagnostic de la
situation}\label{iv.1.2.-diagnostic-de-la-situation}

\textbf{Tableau IV.1 : Analyse SWOT des dispositifs actuels de détection
de fraude au Togo}

{\def\LTcaptype{none} % do not increment counter
\begin{longtable}[]{@{}
  >{\raggedright\arraybackslash}p{(\linewidth - 4\tabcolsep) * \real{0.0667}}
  >{\raggedright\arraybackslash}p{(\linewidth - 4\tabcolsep) * \real{0.4000}}
  >{\raggedright\arraybackslash}p{(\linewidth - 4\tabcolsep) * \real{0.5333}}@{}}
\toprule\noalign{}
\begin{minipage}[b]{\linewidth}\raggedright
\end{minipage} & \begin{minipage}[b]{\linewidth}\raggedright
Forces (S)
\end{minipage} & \begin{minipage}[b]{\linewidth}\raggedright
Faiblesses (W)
\end{minipage} \\
\midrule\noalign{}
\endhead
\bottomrule\noalign{}
\endlastfoot
\textbf{Interne} & S1 — Connaissance fine des clients via les processus
KYC obligatoires & W1 — Règles de détection statiques (seuils fixes,
aucun apprentissage adaptatif) \\
& S2 — Cellules conformité AML déjà structurées dans les grandes banques
& W2 — Faible couverture des transactions mobile money (USSD,
cash-in/out) \\
& S3 — Exigences réglementaires BCEAO imposant des contrôles de base &
W3 — Analyse manuelle des alertes, non scalable au-delà de 500
transactions/jour \\
& S4 — Personnel formé à la conformité et à la gestion des risques & W4
— Délais de détection longs (signalement sous 48h, enquête sous 15
jours) \\
\end{longtable}
}

{\def\LTcaptype{none} % do not increment counter
\begin{longtable}[]{@{}
  >{\raggedright\arraybackslash}p{(\linewidth - 4\tabcolsep) * \real{0.0606}}
  >{\raggedright\arraybackslash}p{(\linewidth - 4\tabcolsep) * \real{0.5455}}
  >{\raggedright\arraybackslash}p{(\linewidth - 4\tabcolsep) * \real{0.3939}}@{}}
\toprule\noalign{}
\begin{minipage}[b]{\linewidth}\raggedright
\end{minipage} & \begin{minipage}[b]{\linewidth}\raggedright
Opportunités (O)
\end{minipage} & \begin{minipage}[b]{\linewidth}\raggedright
Menaces (T)
\end{minipage} \\
\midrule\noalign{}
\endhead
\bottomrule\noalign{}
\endlastfoot
\textbf{Externe} & O1 — Digitalisation rapide du secteur financier
togolais (3,55M utilisateurs mobile money) & T1 — Sophistication
croissante des fraudes (SIM swap, hameçonnage, faux agents) \\
& O2 — Datasets publics de qualité (European Card Transactions 284 807
lignes) & T2 — Multiplication des canaux d’attaque (USSD, mobile money,
virement instantané) \\
& O3 — Outils ML/IA open source matures (XGBoost, SHAP, scikit-learn,
Optuna) & T3 — Ingénierie sociale sur agents et clients, difficile à
détecter automatiquement \\
& O4 — Soutien des régulateurs à l’innovation (SNDC Togo, stratégie
digitale BCEAO) & T4 — Contraintes infrastructurelles (bande passante,
électricité, compétences IT rares) \\
\end{longtable}
}

\emph{Source : Auteur (2025)}

Le diagnostic révèle un contraste marqué entre des fondamentaux solides
(KYC, AML, régulation) et des lacunes opérationnelles importantes. Les
systèmes actuels, essentiellement manuels, ne peuvent suivre le rythme
de croissance des transactions électroniques. Les 3,55 millions
d’utilisateurs de mobile money au Togo (Togo First, 2024) créent une
surface d’attaque croissante que les dispositifs existants ne couvrent
pas.

\subsubsection{IV.1.3. Synthèse diagnostique et besoins
stratégiques}\label{iv.1.3.-synthuxe8se-diagnostique-et-besoins-stratuxe9giques}

La synthèse du diagnostic fait émerger quatre besoins stratégiques
prioritaires :

\textbf{B1 — Automatisation intelligente :} Remplacer les règles
statiques par un système ML capable d’apprendre des patterns de fraude
et de s’adapter. XGBoost démontre sa capacité à atteindre un Recall de
85,02\%, supérieur de 23 points à celui de Random Forest (62\%).

\textbf{B2 — Couverture mobile money :} Étendre la détection aux canaux
USSD, cash-in/out et transferts P2P qui constituent la majorité des
transactions des ménages togolais.

\textbf{B3 — Explicabilité des décisions :} Les analystes fraude et les
auditeurs ont besoin de comprendre le fondement des alertes ML. SHAP
répond à ce besoin en identifiant les variables discriminantes par
transaction.

\textbf{B4 — Sécurité et conformité intégrées :} Le système doit
respecter la Directive BCEAO N°01/2018, la loi N°2020-003 et offrir une
traçabilité complète (logs d’audit, authentification, chiffrement).

\subsubsection{IV.1.4. Vérification des
hypothèses}\label{iv.1.4.-vuxe9rification-des-hypothuxe8ses}

\textbf{HG — Partiellement validée.} XGBoost (Recall=85,02\%, F1=0,23)
surpasse significativement Random Forest (Recall=62\%, F1=0,44) et
Isolation Forest (Recall=14\%, F1=0,18) sur le Recall et l’AUC-PR.
Cependant, le F1 de 0,23 est très inférieur au seuil de confirmation de
0,75, ce qui relativue fortement la portée de la validation.

\textbf{HS1 — Validée.} Le Recall de 85,02\% dépasse le seuil de 85\%
fixé dans le Chapitre II. Les 619 faux négatifs sur 4 127 transactions
frauduleuses représentent 15\% de fraudes non détectées, ce qui
constitue une base acceptable pour un déploiement progressif avec
supervision humaine.

\textbf{HS2 — Vérifiée.} La faisabilité technique de la plateforme
sécurisée est démontrée (authentification JWT HMAC-SHA256, RBAC 3 rôles,
rate limiting, chiffrement TLS 1.3). L’impact sur l’adoption par les
banques togolaises reste à mesurer par une étude terrain.

\textbf{HS3 — Non validée.} Le module SHAP est fonctionnel et permet
d’identifier le top 5 des variables explicatives par alerte via un
endpoint dédié (/explain). Les FP (20,7\%) sont très au-dessus de la
cible de 2\%, ce qui génère un volume d’alertes trop élevé pour une
exploitation opérationnelle sans réduction supplémentaire du seuil.

\textbf{Tableau IV.2 : Synthèse de la vérification des hypothèses}

{\def\LTcaptype{none} % do not increment counter
\begin{longtable}[]{@{}
  >{\raggedright\arraybackslash}p{(\linewidth - 4\tabcolsep) * \real{0.3143}}
  >{\raggedright\arraybackslash}p{(\linewidth - 4\tabcolsep) * \real{0.2571}}
  >{\raggedright\arraybackslash}p{(\linewidth - 4\tabcolsep) * \real{0.4286}}@{}}
\toprule\noalign{}
\begin{minipage}[b]{\linewidth}\raggedright
Hypothèse
\end{minipage} & \begin{minipage}[b]{\linewidth}\raggedright
Verdict
\end{minipage} & \begin{minipage}[b]{\linewidth}\raggedright
Justification
\end{minipage} \\
\midrule\noalign{}
\endhead
\bottomrule\noalign{}
\endlastfoot
HG & Partiellement validée & XGBoost \textgreater{} RF \textgreater{} IF
en Recall, mais F1=0,23 \textless\textless{} seuil 0,75 ; AUC-PR=0,57
\textless{} seuil 0,65 \\
HS1 & Validée & Recall=85,02\% ≥ seuil 85\% ; 619 FN/4 127 fraudes \\
HS2 & Vérifiée & JWT + RBAC 3 rôles + rate limiting fonctionnels ;
impact adoption à mesurer \\
HS3 & Non validée & SHAP fonctionnel ; FP=20,7\%
\textgreater\textgreater{} cible 2\% ; F1=0,23 insuffisant \\
\end{longtable}
}

\emph{Source : Auteur (2025)}

\subsection{IV.2. Intervention proposée et
justification}\label{iv.2.-intervention-proposuxe9e-et-justification}

\subsubsection{IV.2.1. L’intervention : le système
FRAUDX}\label{iv.2.1.-lintervention-le-systuxe8me-fraudx}

L’intervention proposée est le déploiement progressif du système FRAUDX,
une plateforme intégrée de détection de fraude bancaire et mobile money
par intelligence artificielle. FRAUDX combine un moteur de détection
XGBoost optimisé (seuil 0,35), un module d’explicabilité SHAP, un
dashboard interactif avec contrôle d’accès RBAC et un module de feedback
pour l’apprentissage continu.

\subsubsection{IV.2.2. Justification de
l’intervention}\label{iv.2.2.-justification-de-lintervention}

La justification repose sur trois piliers. Sur le plan empirique, les
résultats du Chapitre III démontrent la supériorité de XGBoost
(Recall=85,02\%) sur les approches traditionnelles et les autres modèles
testés, ce qui en fait le candidat principal pour une industrialisation
sous réserve d’amélioration du F1. Sur le plan comparatif, la
littérature africaine récente (Adjovi, 2023 ; Diop \& Ndiaye, 2022 ;
Mensah, 2022) montre que XGBoost est la référence régionale pour la
détection de fraude dans les transactions électroniques. Sur le plan
opérationnel, l’association d’un moteur ML performant, d’un module
d’explicabilité SHAP et d’une sécurité RBAC répond simultanément aux
quatre besoins stratégiques identifiés (automatisation, couverture
mobile money, explicabilité, conformité).

\subsection{IV.3. Objectifs de
l’intervention}\label{iv.3.-objectifs-de-lintervention}

\subsubsection{IV.3.1. Objectif
général}\label{iv.3.1.-objectif-guxe9nuxe9ral}

Déployer un système d’intelligence artificielle opérationnel, sécurisé
et explicable pour la détection en temps réel de la fraude bancaire et
mobile money au Togo, avec pour cibles opérationnelles un Recall ≥ 90\%,
un taux de faux positifs ≤ 3\%, et un temps de réponse inférieur à 100
millisecondes par transaction.

\subsubsection{IV.3.2. Objectifs
spécifiques}\label{iv.3.2.-objectifs-spuxe9cifiques}

\textbf{OSI-1 — Optimisation du modèle :} Adapter et ré-entraîner
XGBoost sur des données locales togolaises avec un objectif de F1 ≥ 0,75
et un Recall ≥ 90\%.

\textbf{OSI-2 — Intégration de l’explicabilité :} Intégrer SHAP dans le
workflow décisionnel des analystes pour que chaque alerte soit
accompagnée de son top 5 des variables explicatives.

\textbf{OSI-3 — Industrialisation de la plateforme :} Déployer la
plateforme sécurisée en production avec authentification JWT, RBAC
complet (3 rôles) et une interface dashboard fonctionnelle.

\textbf{OSI-4 — Renforcement des capacités :} Former au moins 15
analystes fraude et gestionnaires de risques à l’utilisation du système
et à l’interprétation des explications SHAP.

\textbf{OSI-5 — Apprentissage continu :} Mettre en place un pipeline de
feedback permettant de ré-entraîner le modèle mensuellement avec les
transactions validées par les analystes.

\subsection{IV.4. Composantes de l’intervention
envisagée}\label{iv.4.-composantes-de-lintervention-envisaguxe9e}

L’intervention FRAUDX est structurée en sept composantes techniques
interconnectées, couvrant la chaîne complète de la donnée brute à la
décision opérationnelle.

\subsubsection{IV.4.1. Module de collecte et prétraitement des
données}\label{iv.4.1.-module-de-collecte-et-pruxe9traitement-des-donnuxe9es}

Ce module assure l’ingestion en temps réel des flux transactionnels en
provenance des systèmes bancaires et des opérateurs mobile money. Il
normalise les données (encodage des variables catégorielles,
standardisation des montants), applique les règles de pseudonymisation
conformément à la loi N°2020-003 et génère les features engineering
nécessaires au modèle (ratios de fréquence, moyennes glissantes, écarts
à la moyenne client). Le module supporte un volume cible de 50 000
transactions par jour en phase de généralisation.

\subsubsection{IV.4.2. Moteur de détection
XGBoost}\label{iv.4.2.-moteur-de-duxe9tection-xgboost}

Le cœur décisionnel du système repose sur XGBoost, entraîné avec les
hyperparamètres optimisés par Optuna (learning\_rate=0,026,
max\_depth=5, n\_estimators=273, subsample=0,781,
colsample\_bytree=0,803, scale\_pos\_weight=22,4). Le seuil de décision
est fixé à 0,35 (contre 0,5 par défaut) pour maximiser le rappel
(85,02\%) au prix d’un taux de faux positifs de 20,7\%. Ce choix est
délibéré dans un contexte où un faux négatif (fraude non détectée) a un
coût bien supérieur à un faux positif (alerte superflue). Le temps
d’inférence est inférieur à 10 ms par transaction, compatible avec les
exigences du temps réel bancaire.

\subsubsection{IV.4.3. Module d’explicabilité
SHAP}\label{iv.4.3.-module-dexplicabilituxe9-shap}

Chaque alerte générée par XGBoost est accompagnée d’une explication SHAP
locale identifiant le top 5 des variables ayant contribué à la décision,
accessible via un endpoint API dédié (/explain). Par exemple, une
transaction de 250 000 FCFA effectuée à 2h du matin avec un montant 8
fois supérieur à la moyenne du client sera expliquée par les features
“montant”, “hour\_of\_day” et “ratio\_montant\_moyen”. Ces explications
sont visualisées dans le dashboard sous forme de graphiques à barres et
de waterfall plots, permettant à l’analyste de comprendre et de valider
ou infirmer l’alerte en moins de 30 secondes.

\subsubsection{IV.4.4. Dashboard
interactif}\label{iv.4.4.-dashboard-interactif}

Le dashboard FRAUDX est développé avec Streamlit et intègre des cartes
KPI, un graphique d’évolution temporelle, une liste paginée des
transactions avec explications SHAP, un benchmark comparatif et des
visualisations SHAP (importance globale, waterfall plots). L’accès est
sécurisé par authentification JWT.

\subsubsection{IV.4.5. Module de sécurité et
conformité}\label{iv.4.5.-module-de-suxe9curituxe9-et-conformituxe9}

La sécurité est intégrée à tous les niveaux du système. Les données en
transit sont protégées par TLS 1.3, les données au repos par chiffrement
AES-256. L’authentification est gérée par JWT (HMAC-SHA256) avec RBAC à
trois rôles (analyste, superviseur, administrateur). Chaque action
(prédiction, qualification d’alerte, modification de seuil) est
horodatée et conservée dans des logs d’audit pendant 10 ans,
conformément à la réglementation BCEAO.

\subsubsection{IV.4.6. Module de feedback et apprentissage
continu}\label{iv.4.6.-module-de-feedback-et-apprentissage-continu}

Les qualifications effectuées par les analystes (fraude confirmée / faux
positif) sont stockées dans une base dédiée. Un pipeline automatisé
ré-entraîne XGBoost mensuellement en intégrant ces nouvelles données
labellisées, permettant au modèle de s’adapter aux nouvelles typologies
de fraude. Ce mécanisme réduit progressivement le taux de faux positifs
(cible ≤ 3\% à 6 mois) et maintient un Recall élevé.

\subsubsection{IV.4.7. API d’intégration et
connecteurs}\label{iv.4.7.-api-dintuxe9gration-et-connecteurs}

Une API RESTful sécurisée expose les fonctionnalités du système aux
systèmes bancaires existants (core banking, switch mobile money,
plateformes AML). Les connecteurs standardisés (format JSON,
authentification JWT avec RBAC) permettent une intégration rapide sans
modification des systèmes d’information des banques partenaires. Le
temps de réponse API est inférieur à 100 ms pour 95\% des requêtes.

\subsection{IV.5. Stratégies d’action et
périmètre}\label{iv.5.-stratuxe9gies-daction-et-puxe9rimuxe8tre}

\subsubsection{IV.5.1. Stratégies
d’action}\label{iv.5.1.-stratuxe9gies-daction}

\textbf{Phase pilote — Mois 1 à 6 :} Déploiement chez une banque
partenaire (SUNU Bank Togo) avec les flux de transactions bancaires
classiques (CB, virement, paiement). Montée en charge progressive de 1
000 à 10 000 transactions par jour. Cette phase valide l’intégration
technique, la performance du modèle sur des données réelles et
l’acceptation par les analystes. Les objectifs sont un Recall ≥ 85\% et
un taux de FP ≤ 5\%.

\textbf{Extension mobile money — Mois 7 à 12 :} Intégration des flux
mobile money (USSD, cash-in/out, transferts P2P). Ajout de features
spécifiques (type d’opérateur, heure d’envoi, fréquence hebdomadaire).
Montée à 50 000 transactions par jour. L’objectif est d’atteindre un
Recall ≥ 90\% et un FP ≤ 3\% sur l’ensemble des canaux.

\textbf{Généralisation — Mois 13 à 24 :} Extension à 3 à 5 banques et
opérateurs mobile money. Mise en place d’un centre de veille mutualisé.
Création d’un consortium de partage de données anonymisées. L’objectif
est de couvrir 60\% des transactions électroniques du Togo.

\subsubsection{IV.5.2. Périmètre de
l’intervention}\label{iv.5.2.-puxe9rimuxe8tre-de-lintervention}

Le périmètre fonctionnel couvre trois canaux de transaction (bancaire
classique, mobile money USSD, transferts P2P) et trois types d’acteurs
(banques, opérateurs mobile money, régulateurs). Le périmètre
géographique se limite au Togo dans un premier temps, avec une
perspective d’extension à l’UEMOA. Sont exclus du périmètre initial la
détection de fraude interne (collusion employé-client), la fraude
documentaire et le blanchiment d’argent (AML), ces derniers relevant de
dispositifs spécialisés existants.

\subsection{IV.6. Étude de
faisabilité}\label{iv.6.-uxe9tude-de-faisabilituxe9}

\subsubsection{IV.6.1. Faisabilité
économique}\label{iv.6.1.-faisabilituxe9-uxe9conomique}

Le budget d’investissement et d’exploitation est estimé sur trois ans en
tenant compte des coûts d’infrastructure on-premise (option privilégiée
pour la souveraineté des données), de développement logiciel, de
formation et de maintenance.

\textbf{Tableau IV.3 : Budget détaillé de l’intervention FRAUDX}

{\def\LTcaptype{none} % do not increment counter
\begin{longtable}[]{@{}
  >{\raggedright\arraybackslash}p{(\linewidth - 8\tabcolsep) * \real{0.3103}}
  >{\raggedright\arraybackslash}p{(\linewidth - 8\tabcolsep) * \real{0.1552}}
  >{\raggedright\arraybackslash}p{(\linewidth - 8\tabcolsep) * \real{0.1552}}
  >{\raggedright\arraybackslash}p{(\linewidth - 8\tabcolsep) * \real{0.1552}}
  >{\raggedright\arraybackslash}p{(\linewidth - 8\tabcolsep) * \real{0.2241}}@{}}
\toprule\noalign{}
\begin{minipage}[b]{\linewidth}\raggedright
Poste de dépense
\end{minipage} & \begin{minipage}[b]{\linewidth}\raggedright
Année 1
\end{minipage} & \begin{minipage}[b]{\linewidth}\raggedright
Année 2
\end{minipage} & \begin{minipage}[b]{\linewidth}\raggedright
Année 3
\end{minipage} & \begin{minipage}[b]{\linewidth}\raggedright
Total 3 ans
\end{minipage} \\
\midrule\noalign{}
\endhead
\bottomrule\noalign{}
\endlastfoot
Infrastructure serveur (ML + API + BDD + stockage) & 12 500 € & 2 000 €
& 2 000 € & 16 500 € \\
Développement et optimisation ML & 30 000 € & 8 000 € & 8 000 € & 46 000
€ \\
Développement dashboard et interface & 15 000 € & 3 000 € & 3 000 € & 21
000 € \\
Formation (15 analystes, 5 jours) & 10 000 € & 3 000 € & 3 000 € & 16
000 € \\
Maintenance évolutive et corrective & 5 000 € & 8 000 € & 10 000 € & 23
000 € \\
\textbf{Total} & \textbf{72 500 €} & \textbf{24 000 €} & \textbf{26 000
€} & \textbf{122 500 €} \\
\end{longtable}
}

\emph{Source : Auteur (2025)}

\textbf{Calcul du retour sur investissement (ROI) :}

Les pertes annuelles estimées pour une banque de taille moyenne au Togo
sont de 300 000 € (fraudes non détectées, remboursements clients,
atteinte à la réputation). Sur la base des résultats de XGBoost
(Recall=85\%), une réduction conservatrice de 40\% des pertes est
attendue, soit une économie annuelle de 120 000 €.

\begin{itemize}
\tightlist
\item
  Économie totale sur 3 ans : 120 000 € × 3 = 360 000 €
\item
  Coût total sur 3 ans : 122 500 €
\item
  ROI = (360 000 - 122 500) / 122 500 = \textbf{194\%}
\item
  Seuil de rentabilité : 122 500 / 120 000 ≈ 12,2 mois
\end{itemize}

Le système devient donc rentable dès la fin de la première année
d’exploitation.

\subsubsection{IV.6.2. Faisabilité
sociale}\label{iv.6.2.-faisabilituxe9-sociale}

L’acceptation par les utilisateurs finaux est un facteur critique.
FRAUDX est conçu comme un outil d’aide à la décision : l’analyste
conserve le pouvoir de validation finale et peut infirmer une alerte
générée par le modèle. Les mesures d’accompagnement incluent une
formation obligatoire de 5 jours (analyse des alertes SHAP, navigation
dashboard, qualification des transactions), une phase de transition de 3
mois en mode parallèle (sans impact sur les opérations existantes), et
un système de feedback continu permettant aux utilisateurs de signaler
les dysfonctionnements ou les améliorations souhaitées. Le risque de
rejet technologique est atténué par l’accent mis sur l’explicabilité
(SHAP) et le maintien de l’autonomie décisionnelle des analystes.

\subsubsection{IV.6.3. Faisabilité
technique}\label{iv.6.3.-faisabilituxe9-technique}

L’infrastructure requise pour le déploiement on-premise comprend : un
serveur ML (32 vCPU, 64 Go RAM, GPU optionnel — 4 000 €), un serveur API
(8 vCPU, 32 Go RAM — 2 000 €), un serveur de base de données (16 vCPU,
64 Go RAM, SSD — 3 000 €), un NAS 10 To pour les sauvegardes et logs
d’audit (1 500 €), et une solution de sécurité (firewall applicatif,
WAF, VPN — 2 000 €). Soit un total de 12 500 € en investissement
initial. La pile logicielle utilise exclusivement des technologies open
source (Python, XGBoost, SHAP, Dash, PostgreSQL) garantissant l’absence
de coûts de licence. Le temps d’inférence mesuré (\textless{} 10 ms par
transaction) est compatible avec les contraintes du temps réel bancaire.

\subsubsection{IV.6.4. Faisabilité
environnementale}\label{iv.6.4.-faisabilituxe9-environnementale}

L’empreinte environnementale du système est limitée. L’infrastructure
on-premise (4 serveurs) consomme environ 3 500 kWh par an, soit
l’équivalent de 1,2 tonne de CO₂ (estimation ADEME, mix électrique
ouest-africain). L’optimisation des hyperparamètres par Optuna a permis
de réduire la taille du modèle (100 arbres, profondeur 5) sans sacrifice
de performance. L’architecture modulaire permet une mise à l’échelle
progressive, évitant le surdimensionnement initial. Aucun matériel
spécifique à forte empreinte carbone (GPU dédié) n’est requis pour
l’inférence.

\subsection{IV.7. Perspectives et
limites}\label{iv.7.-perspectives-et-limites}

Le déploiement de FRAUDX ouvre plusieurs perspectives. À court terme,
l’établissement d’un partenariat avec une banque ou un opérateur mobile
money togolais permettrait une validation sur données réelles. À moyen
terme, l’extension à l’échelle de l’UEMOA (huit pays, même cadre
réglementaire BCEAO) multiplierait l’impact. À long terme, l’intégration
de techniques avancées comme l’apprentissage fédéré (préservation de la
confidentialité des données) et le deep learning (LSTM pour les
séquences temporelles, Transformers pour les patterns complexes)
constituerait une évolution naturelle.

Cependant, plusieurs limites doivent être reconnues. L’absence de
validation sur des données togolaises réelles limite la généralisabilité
des résultats. L’approche qualitative (entretiens, questionnaire TAM)
n’a pu être réalisée, ce qui restreint la compréhension des freins à
l’adoption. Le LSTM n’a pas été implémenté en raison des contraintes de
calcul GPU. Les coûts estimés sont préliminaires et sujets à révision.
Enfin, le taux de faux positifs (4,77\%) reste supérieur à la cible de
2\%, ce qui pourrait générer une fatigue des alertes en production.

\subsection{Conclusion du chapitre}\label{conclusion-du-chapitre-3}

Ce quatrième chapitre a établi un diagnostic approfondi de la situation
de la détection de fraude au Togo, confirmant la pertinence d’une
intervention fondée sur l’intelligence artificielle. La vérification des
hypothèses montre que HS1 est validée, HG partiellement validée, HS2
vérifiée, et HS3 non validée. L’intervention FRAUDX, détaillée en sept
composantes techniques interconnectées, est justifiée par les
performances empiriques de XGBoost et les besoins stratégiques
identifiés. L’étude de faisabilité confirme la viabilité économique (ROI
de 194\% sur trois ans, rentabilité à 12,2 mois), sociale (outil d’aide
à la décision, formation accompagnée) et technique (technologies open
source, temps d’inférence \textless{} 10 ms). Les perspectives
d’extension à l’UEMOA et d’intégration de techniques avancées ouvrent la
voie à une industrialisation régionale du système.

\begin{center}\rule{0.5\linewidth}{0.5pt}\end{center}

\section{CONCLUSION GENERALE}\label{conclusion-generale}

Ce mémoire avait pour objectif de concevoir un système d’intelligence
artificielle performant et sécurisé pour la détection de la fraude
bancaire dans le contexte spécifique du Togo.

Le Chapitre I a posé les fondements théoriques, montrant que la fraude
bancaire au Togo présente des caractéristiques spécifiques liées à la
prédominance du mobile money. Le Chapitre II a défini la méthodologie
quantitative comparative. Le Chapitre III a présenté les résultats
expérimentaux montrant la supériorité relative de XGBoost
(Recall=85,02\%, F1=0,23, AUC-PR=0,57) et démontré la faisabilité
technique via la PoC FRAUDX. Le Chapitre IV a établi un diagnostic
approfondi (SWOT, quatre besoins stratégiques, vérification des
hypothèses) et proposé une intervention structurée en sept composantes
techniques avec un déploiement progressif sur 24 mois et un ROI estimé à
194\% (rentabilité à 12,2 mois).

\textbf{Vérification des hypothèses :}

{\def\LTcaptype{none} % do not increment counter
\begin{longtable}[]{@{}
  >{\raggedright\arraybackslash}p{(\linewidth - 2\tabcolsep) * \real{0.5500}}
  >{\raggedright\arraybackslash}p{(\linewidth - 2\tabcolsep) * \real{0.4500}}@{}}
\toprule\noalign{}
\begin{minipage}[b]{\linewidth}\raggedright
Hypothèse
\end{minipage} & \begin{minipage}[b]{\linewidth}\raggedright
Verdict
\end{minipage} \\
\midrule\noalign{}
\endhead
\bottomrule\noalign{}
\endlastfoot
HG — XGBoost surpasse RF et IF & \textbf{Partiellement validée} \\
HS1 — Le ML réduit les faux négatifs & \textbf{Validée} \\
HS2 — Faisabilité technique de la plateforme RBAC & \textbf{Vérifiée} \\
HS3 — L’explicabilité SHAP facilite le contrôle des FP & \textbf{Non
validée} \\
\end{longtable}
}

Les principales limites sont l’absence de données locales réelles, la
non-réalisation des entretiens qualitatifs, la non-implantation du LSTM
(contraintes GPU), le taux de FP (20,7\%) très supérieur à la cible de
2\%, et le périmètre limité au Togo.

Les perspectives incluent un partenariat avec une banque togolaise pour
des données réelles, l’extension à l’UEMOA, l’apprentissage fédéré, le
deep learning, et l’étude d’acceptabilité par questionnaire TAM et
entretiens.

En définitive, cette étude apporte une première démonstration de la
faisabilité d’un système d’IA fondé sur XGBoost, SHAP et RBAC pour la
détection de fraude bancaire en contexte togolais. Les résultats
obtenus, bien que partiels, constituent une base solide pour un
déploiement pilote et des travaux futurs.

\begin{center}\rule{0.5\linewidth}{0.5pt}\end{center}

\section{BIBLIOGRAPHIE ET
WEBOGRAPHIE}\label{bibliographie-et-webographie}

\subsection{Bibliographie}\label{bibliographie}

Akiba, T., Sano, S., Yanase, T., Ohta, T., \& Koyama, M. (2019). Optuna:
A next-generation hyperparameter optimization framework.
\emph{Proceedings of the 25th ACM SIGKDD International Conference on
Knowledge Discovery and Data Mining}, 2623-2631.

Bhattacharyya, S., Jha, S., Tharakunnel, K., \& Westland, J. C. (2011).
Data mining for credit card fraud: A comparative study. \emph{Decision
Support Systems}, 50(3), 602-613.

Bolton, R. J., \& Hand, D. J. (2002). Statistical fraud detection: A
review. \emph{Statistical Science}, 17(3), 235-255.

Breiman, L. (2001). Random forests. \emph{Machine Learning}, 45(1),
5-32.

Chawla, N. V., Bowyer, K. W., Hall, L. O., \& Kegelmeyer, W. P. (2002).
SMOTE: Synthetic minority over-sampling technique. \emph{Journal of
Artificial Intelligence Research}, 16, 321-357.

Chen, T., \& Guestrin, C. (2016). XGBoost: A scalable tree boosting
system. \emph{Proceedings of the 22nd ACM SIGKDD International
Conference on Knowledge Discovery and Data Mining}, 785-794.

Chen, Y., Zhang, W., \& Liu, H. (2026). SAGE: A multi-agent LLM
framework for interpretable fraud detection. \emph{arXiv preprint
arXiv:2606.08146}.

Dal Pozzolo, A., Caelen, O., Le Borgne, Y.-A., Waterschoot, S., \&
Bontempi, G. (2014). Learned lessons in credit card fraud detection from
a practitioner perspective. \emph{Expert Systems with Applications},
41(10), 4915-4928.

Davis, F. D. (1989). Perceived usefulness, perceived ease of use, and
user acceptance of information technology. \emph{MIS Quarterly}, 13(3),
319-340.

Hochreiter, S., \& Schmidhuber, J. (1997). Long short-term memory.
\emph{Neural Computation}, 9(8), 1735-1780.

Liu, F. T., Ting, K. M., \& Zhou, Z.-H. (2008). Isolation forest.
\emph{2008 Eighth IEEE International Conference on Data Mining},
413-422.

Liu, F. T., Ting, K. M., \& Zhou, Z.-H. (2012). Isolation-based anomaly
detection. \emph{ACM Transactions on Knowledge Discovery from Data},
6(1), 1-39.

Lundberg, S. M., \& Lee, S.-I. (2017). A unified approach to
interpreting model predictions. \emph{Advances in Neural Information
Processing Systems}, 30, 4765-4774.

Lundberg, S. M., Erion, G., Chen, H., DeGrave, A., Prutkin, J. M., Nair,
B., … \& Lee, S.-I. (2020). From local explanations to global
understanding with explainable AI for trees. \emph{Nature Machine
Intelligence}, 2(1), 56-67.

Samuel, A. L. (1959). Some studies in machine learning using the game of
checkers. \emph{IBM Journal of Research and Development}, 3(3), 210-229.

Shapley, L. S. (1953). A value for n-person games. In H. W. Kuhn \& A.
W. Tucker (Eds.), \emph{Contributions to the Theory of Games} (Vol. 2,
pp.~307-317). Princeton University Press.

Adjovi, E. (2023). Détection de fraude mobile money par régression
logistique au Bénin. \emph{Revue de l’Innovation et de la Technologie},
5(3), 78-91.

Diop, M., \& Ndiaye, S. (2022). Amélioration de la détection de fraude
bancaire par XGBoost au Sénégal. \emph{Annales de l’Université Cheikh
Anta Diop}, 28(1), 112-128.

Kouamé, A. K. (2021). Détection de fraude bancaire par apprentissage
automatique en Côte d’Ivoire. \emph{Revue Africaine de Recherche en
Informatique}, 14(2), 45-62.

Mensah, K. (2022). Mobile money fraud detection using XGBoost and SMOTE
in Ghana. \emph{West African Journal of Applied Computing}, 9(1), 34-51.

Okonkwo, C., Eze, P., \& Okafor, N. (2020). Ensemble learning for fraud
detection in Nigerian banking sector. \emph{Journal of African Fintech},
3(2), 156-173.

Qian, Z., Wang, L., \& Li, J. (2025). FraudGuess: Explainable fraud
pattern discovery via micro-clustering and visual analytics. \emph{arXiv
preprint arXiv:2509.15493}.

Quivy, R., \& Van Campenhoudt, L. (2006). \emph{Manuel de recherche en
sciences sociales} (3e éd.). Dunod.

\subsection{Webographie}\label{webographie}

ANCY. (2025a, février). \emph{Alerte aux arnaques par faux transfert
Mobile Money}. Agence Nationale de Cybersécurité, République Togolaise.

ANCY. (2025b, mars). \emph{Alerte à l’usurpation d’agent Mobile Money et
à la fausse réidentification}. Agence Nationale de Cybersécurité,
République Togolaise.

BCEAO. (2024). \emph{Rapport sur les systèmes de paiement dans l’UEMOA}.

CERT-TG. (2025, avril). \emph{Alerte aux plateformes frauduleuses de
vente et d’investissement}. Centre Togolais de Réponse aux Incidents de
Sécurité Informatique.

République Togolaise. (2020). \emph{Loi N°2020-003 du 20 février 2020
relative à la protection des données à caractère personnel}.

Togo First. (2024, novembre). Mobile money au Togo : 3,55 millions
d’utilisateurs, l’ARCEP dresse le portrait-robot. \emph{Togo First}.

UEMOA. (2018). \emph{Directive N°01/2018/CM/UEMOA relative aux systèmes
de paiement}.

UEMOA. (2020). \emph{Règlement N°01/2020/CM/UEMOA sur les services de
paiement mobile}.

\begin{center}\rule{0.5\linewidth}{0.5pt}\end{center}

\section{ANNEXES}\label{annexes}

\subsection{Annexe A : Grille d’entretien
semi-directif}\label{annexe-a-grille-dentretien-semi-directif}

\textbf{Profils cibles :} DSI/IT, Conformité/KYC-AML, Gestion des
risques / Analyste fraude \textbf{Durée estimée :} 25-30 minutes

\textbf{Thème 1 — Typologies de fraude observées :} 1. Quels types de
fraudes sont les plus fréquents dans votre institution ? 2.
Observez-vous des fraudes spécifiques au mobile money (SIM swap, USSD,
ingénierie sociale) ? 3. Comment ces fraudes ont-elles évolué ces 3-5
dernières années ?

\textbf{Thème 2 — Systèmes de détection actuels :} 4. Quels outils et
méthodes utilisez-vous actuellement ? 5. Quelles sont les principales
limites de ces systèmes ? 6. Un système de détection en temps réel
serait-il utile ?

\textbf{Thème 3 — IA et Machine Learning :} 7. Quel est votre niveau de
familiarité avec l’IA/ML ? 8. L’IA pourrait-elle améliorer la détection
? 9. Quels seraient les freins à l’adoption de l’IA dans une banque
togolaise ? 10. Un modèle qui explique ses décisions serait-il mieux
adopté ?

\textbf{Thème 4 — Contraintes et besoins :} 11. Quelles contraintes
infrastructurelles pourraient limiter le déploiement ? 12. La
disponibilité des données est-elle un problème ? 13. Quelles sont vos
attentes en matière de conformité réglementaire ?

\subsection{Annexe B : Guide TAM complet — Adoption du système
FRAUDX}\label{annexe-b-guide-tam-complet-adoption-du-systuxe8me-fraudx}

\subsubsection{Partie 1 — Cadre théorique du Technology Acceptance Model
(TAM)}\label{partie-1-cadre-thuxe9orique-du-technology-acceptance-model-tam}

Le Technology Acceptance Model (Davis, 1989) est le cadre théorique le
plus éprouvé pour prédire et expliquer l’adoption d’une technologie par
ses utilisateurs. Il repose sur deux construits fondamentaux :

{\def\LTcaptype{none} % do not increment counter
\begin{longtable}[]{@{}
  >{\raggedright\arraybackslash}p{(\linewidth - 4\tabcolsep) * \real{0.2500}}
  >{\raggedright\arraybackslash}p{(\linewidth - 4\tabcolsep) * \real{0.2500}}
  >{\raggedright\arraybackslash}p{(\linewidth - 4\tabcolsep) * \real{0.5000}}@{}}
\toprule\noalign{}
\begin{minipage}[b]{\linewidth}\raggedright
Construit
\end{minipage} & \begin{minipage}[b]{\linewidth}\raggedright
Définition
\end{minipage} & \begin{minipage}[b]{\linewidth}\raggedright
Application à FRAUDX
\end{minipage} \\
\midrule\noalign{}
\endhead
\bottomrule\noalign{}
\endlastfoot
\textbf{Utilité perçue (PU)} & Degré auquel un utilisateur croit que la
technologie améliorera sa performance & Le système détecte-t-il plus de
fraudes ? Réduit-il le temps d’analyse ? \\
\textbf{Facilité d’utilisation perçue (PEOU)} & Degré auquel
l’utilisateur croit que l’usage sera sans effort & Le dashboard est-il
intuitif ? Les explications SHAP sont-elles compréhensibles ? \\
\textbf{Confiance (Trust)} & Extension postérieure (Gefen et al., 2003)
: croyance en la fiabilité du système & L’utilisateur accepte-t-il les
décisions du ML ? L’explicabilité renforce-t-elle la confiance ? \\
\textbf{Intention d’adoption (BI)} & Variable dépendante ultime :
l’intention d’utiliser le système & L’utilisateur recommanderait-il
FRAUDX ? Participerait-il à un pilote ? \\
\end{longtable}
}

Ces deux construits influencent l’intention d’adoption (BI), qui
elle-même prédit l’usage réel. Dans ce mémoire, deux variables
contextuelles sont ajoutées : la confiance (TR) et les facteurs
contextuels locaux (FC — infrastructure, données, réglementation).

\subsubsection{Partie 2 — Guide d’entretien semi-structuré
TAM}\label{partie-2-guide-dentretien-semi-structuruxe9-tam}

\textbf{Profils cibles :} Analystes fraude, Gestionnaires de risques,
Conformité/KYC-AML, DSI, Responsables mobile money \textbf{Durée estimée
:} 30-40 minutes \textbf{Modalité :} Entretien individuel, semi-directif
(guide souple, rebond possible)

\paragraph{Thème 1 — Perception du problème et besoins
(Contextualisation)}\label{thuxe8me-1-perception-du-probluxe8me-et-besoins-contextualisation}

{\def\LTcaptype{none} % do not increment counter
\begin{longtable}[]{@{}
  >{\raggedright\arraybackslash}p{(\linewidth - 6\tabcolsep) * \real{0.1026}}
  >{\raggedright\arraybackslash}p{(\linewidth - 6\tabcolsep) * \real{0.2564}}
  >{\raggedright\arraybackslash}p{(\linewidth - 6\tabcolsep) * \real{0.3846}}
  >{\raggedright\arraybackslash}p{(\linewidth - 6\tabcolsep) * \real{0.2564}}@{}}
\toprule\noalign{}
\begin{minipage}[b]{\linewidth}\raggedright
N°
\end{minipage} & \begin{minipage}[b]{\linewidth}\raggedright
Question
\end{minipage} & \begin{minipage}[b]{\linewidth}\raggedright
Dimension TAM
\end{minipage} & \begin{minipage}[b]{\linewidth}\raggedright
Objectif
\end{minipage} \\
\midrule\noalign{}
\endhead
\bottomrule\noalign{}
\endlastfoot
1 & Comment décririez-vous le problème de la fraude bancaire dans votre
institution aujourd’hui ? & Contexte & Établir le référentiel partagé \\
2 & Quelles sont, selon vous, les lacunes des systèmes actuels de
détection ? & Contexte & Identifier les irritants \\
3 & Qu’attendez-vous prioritairement d’un nouveau système de détection ?
& PU (besoin) & Faire émerger les critères implicites \\
\end{longtable}
}

\paragraph{Thème 2 — Utilité perçue du système FRAUDX
(PU)}\label{thuxe8me-2-utilituxe9-peruxe7ue-du-systuxe8me-fraudx-pu}

{\def\LTcaptype{none} % do not increment counter
\begin{longtable}[]{@{}
  >{\raggedright\arraybackslash}p{(\linewidth - 6\tabcolsep) * \real{0.0976}}
  >{\raggedright\arraybackslash}p{(\linewidth - 6\tabcolsep) * \real{0.2439}}
  >{\raggedright\arraybackslash}p{(\linewidth - 6\tabcolsep) * \real{0.3659}}
  >{\raggedright\arraybackslash}p{(\linewidth - 6\tabcolsep) * \real{0.2927}}@{}}
\toprule\noalign{}
\begin{minipage}[b]{\linewidth}\raggedright
N°
\end{minipage} & \begin{minipage}[b]{\linewidth}\raggedright
Question
\end{minipage} & \begin{minipage}[b]{\linewidth}\raggedright
Dimension TAM
\end{minipage} & \begin{minipage}[b]{\linewidth}\raggedright
Indicateur
\end{minipage} \\
\midrule\noalign{}
\endhead
\bottomrule\noalign{}
\endlastfoot
4 & Si FRAUDX vous signalait une transaction suspecte avec un score de
risque et une explication SHAP, en quoi cela changerait-il votre analyse
? & PU — Performance & Gain perçu vs.~méthode actuelle \\
5 & Pensez-vous qu’un tel système pourrait détecter des fraudes qui vous
échappent aujourd’hui ? Si oui, lesquelles ? & PU — Efficacité &
Coverage des typologies \\
6 & Le temps de traitement des alertes vous semble-t-il un levier
d’amélioration important ? & PU — Productivité & Réduction du temps
perçu \\
\end{longtable}
}

\paragraph{Thème 3 — Facilité d’utilisation perçue
(PEOU)}\label{thuxe8me-3-facilituxe9-dutilisation-peruxe7ue-peou}

{\def\LTcaptype{none} % do not increment counter
\begin{longtable}[]{@{}
  >{\raggedright\arraybackslash}p{(\linewidth - 6\tabcolsep) * \real{0.0976}}
  >{\raggedright\arraybackslash}p{(\linewidth - 6\tabcolsep) * \real{0.2439}}
  >{\raggedright\arraybackslash}p{(\linewidth - 6\tabcolsep) * \real{0.3659}}
  >{\raggedright\arraybackslash}p{(\linewidth - 6\tabcolsep) * \real{0.2927}}@{}}
\toprule\noalign{}
\begin{minipage}[b]{\linewidth}\raggedright
N°
\end{minipage} & \begin{minipage}[b]{\linewidth}\raggedright
Question
\end{minipage} & \begin{minipage}[b]{\linewidth}\raggedright
Dimension TAM
\end{minipage} & \begin{minipage}[b]{\linewidth}\raggedright
Indicateur
\end{minipage} \\
\midrule\noalign{}
\endhead
\bottomrule\noalign{}
\endlastfoot
7 & En regardant ces deux visuels (bar plot SHAP, waterfall plot), que
comprenez-vous de la décision du modèle ? & PEOU — Compréhension &
Clarté perçue \\
8 & Combien de temps de formation pensez-vous qu’un analyste
non-spécialiste en IA aurait besoin pour utiliser ce tableau de bord ? &
PEOU — Apprentissage & Effort perçu \\
9 & Qu’est-ce qui pourrait rendre l’interface difficile à utiliser pour
vos équipes ? & PEOU — Obstacles & Freins UX \\
\end{longtable}
}

\paragraph{Thème 4 — Confiance et explicabilité
(TR)}\label{thuxe8me-4-confiance-et-explicabilituxe9-tr}

{\def\LTcaptype{none} % do not increment counter
\begin{longtable}[]{@{}
  >{\raggedright\arraybackslash}p{(\linewidth - 6\tabcolsep) * \real{0.0976}}
  >{\raggedright\arraybackslash}p{(\linewidth - 6\tabcolsep) * \real{0.2439}}
  >{\raggedright\arraybackslash}p{(\linewidth - 6\tabcolsep) * \real{0.3659}}
  >{\raggedright\arraybackslash}p{(\linewidth - 6\tabcolsep) * \real{0.2927}}@{}}
\toprule\noalign{}
\begin{minipage}[b]{\linewidth}\raggedright
N°
\end{minipage} & \begin{minipage}[b]{\linewidth}\raggedright
Question
\end{minipage} & \begin{minipage}[b]{\linewidth}\raggedright
Dimension TAM
\end{minipage} & \begin{minipage}[b]{\linewidth}\raggedright
Indicateur
\end{minipage} \\
\midrule\noalign{}
\endhead
\bottomrule\noalign{}
\endlastfoot
10 & Recevoir une explication (ex. : « Cette transaction est suspecte
car le montant est 3× supérieur à la moyenne et l’IP est inhabituelle »)
vous inciterait-elle à faire davantage confiance à l’alerte ? & TR —
Explicabilité & Valeur ajoutée du SHAP \\
11 & À partir de quel niveau de confiance dans le système
accepteriez-vous de ne pas vérifier manuellement chaque alerte ? & TR —
Délégation & Seuil d’autonomie \\
12 & Qu’est-ce qui vous freinerait dans l’adoption d’une décision prise
par une IA ? & TR — Réticence & Facteurs de défiance \\
\end{longtable}
}

\paragraph{Thème 5 — Adoption et passage à l’échelle
(BI)}\label{thuxe8me-5-adoption-et-passage-uxe0-luxe9chelle-bi}

{\def\LTcaptype{none} % do not increment counter
\begin{longtable}[]{@{}
  >{\raggedright\arraybackslash}p{(\linewidth - 6\tabcolsep) * \real{0.0976}}
  >{\raggedright\arraybackslash}p{(\linewidth - 6\tabcolsep) * \real{0.2439}}
  >{\raggedright\arraybackslash}p{(\linewidth - 6\tabcolsep) * \real{0.3659}}
  >{\raggedright\arraybackslash}p{(\linewidth - 6\tabcolsep) * \real{0.2927}}@{}}
\toprule\noalign{}
\begin{minipage}[b]{\linewidth}\raggedright
N°
\end{minipage} & \begin{minipage}[b]{\linewidth}\raggedright
Question
\end{minipage} & \begin{minipage}[b]{\linewidth}\raggedright
Dimension TAM
\end{minipage} & \begin{minipage}[b]{\linewidth}\raggedright
Indicateur
\end{minipage} \\
\midrule\noalign{}
\endhead
\bottomrule\noalign{}
\endlastfoot
13 & Seriez-vous prêt à recommander FRAUDX à votre hiérarchie pour un
projet pilote ? & BI — Intention & Adoption \\
14 & Selon vous, quels seraient les trois principaux obstacles au
déploiement d’un tel système au Togo ? & FC — Contraintes & Freins
contextuels \\
15 & Quelles adaptations locales seraient nécessaires pour que le
système soit pleinement pertinent au Togo ? & FC — Adéquation &
Spécificités locales \\
\end{longtable}
}

\subsubsection{Partie 3 — Questionnaire quantitatif TAM
(Auto-administré)}\label{partie-3-questionnaire-quantitatif-tam-auto-administruxe9}

\textbf{Échelle :} Likert à 5 niveaux (1 = Pas du tout d’accord, 5 =
Tout à fait d’accord)

\paragraph{Section A — Profil du
répondant}\label{section-a-profil-du-ruxe9pondant}

A1. Catégorie de poste : Analyste fraude / Gestionnaire risques /
Conformité / DSI / Mobile money A2. Années d’expérience dans le secteur
financier : \textless{} 2 ans / 2-5 ans / 5-10 ans / \textgreater{} 10
ans A3. Utilisez-vous actuellement un système automatisé de détection de
fraude ? Oui / Non / Partiellement

\paragraph{Section B — Utilité perçue
(PU)}\label{section-b-utilituxe9-peruxe7ue-pu}

{\def\LTcaptype{none} % do not increment counter
\begin{longtable}[]{@{}
  >{\raggedright\arraybackslash}p{(\linewidth - 12\tabcolsep) * \real{0.1622}}
  >{\raggedright\arraybackslash}p{(\linewidth - 12\tabcolsep) * \real{0.1622}}
  >{\raggedright\arraybackslash}p{(\linewidth - 12\tabcolsep) * \real{0.1351}}
  >{\raggedright\arraybackslash}p{(\linewidth - 12\tabcolsep) * \real{0.1351}}
  >{\raggedright\arraybackslash}p{(\linewidth - 12\tabcolsep) * \real{0.1351}}
  >{\raggedright\arraybackslash}p{(\linewidth - 12\tabcolsep) * \real{0.1351}}
  >{\raggedright\arraybackslash}p{(\linewidth - 12\tabcolsep) * \real{0.1351}}@{}}
\toprule\noalign{}
\begin{minipage}[b]{\linewidth}\raggedright
Code
\end{minipage} & \begin{minipage}[b]{\linewidth}\raggedright
Item
\end{minipage} & \begin{minipage}[b]{\linewidth}\raggedright
PU1
\end{minipage} & \begin{minipage}[b]{\linewidth}\raggedright
PU2
\end{minipage} & \begin{minipage}[b]{\linewidth}\raggedright
PU3
\end{minipage} & \begin{minipage}[b]{\linewidth}\raggedright
PU4
\end{minipage} & \begin{minipage}[b]{\linewidth}\raggedright
PU5
\end{minipage} \\
\midrule\noalign{}
\endhead
\bottomrule\noalign{}
\endlastfoot
PU1 & FRAUDX améliorerait significativement la détection des fraudes
dans mon institution & 1 & 2 & 3 & 4 & 5 \\
PU2 & FRAUDX permettrait de détecter des fraudes invisibles aux méthodes
actuelles & 1 & 2 & 3 & 4 & 5 \\
PU3 & FRAUDX réduirait le temps d’analyse des alertes & 1 & 2 & 3 & 4 &
5 \\
PU4 & FRAUDX aiderait à prioriser les alertes critiques & 1 & 2 & 3 & 4
& 5 \\
PU5 & Le système serait utile pour le reporting réglementaire & 1 & 2 &
3 & 4 & 5 \\
\end{longtable}
}

Score PU = moyenne de PU1 à PU5 (cible ≥ 3,5/5 pour validation)

\paragraph{Section C — Facilité d’utilisation perçue
(PEOU)}\label{section-c-facilituxe9-dutilisation-peruxe7ue-peou}

{\def\LTcaptype{none} % do not increment counter
\begin{longtable}[]{@{}
  >{\raggedright\arraybackslash}p{(\linewidth - 12\tabcolsep) * \real{0.2222}}
  >{\raggedright\arraybackslash}p{(\linewidth - 12\tabcolsep) * \real{0.2222}}
  >{\raggedright\arraybackslash}p{(\linewidth - 12\tabcolsep) * \real{0.1111}}
  >{\raggedright\arraybackslash}p{(\linewidth - 12\tabcolsep) * \real{0.1111}}
  >{\raggedright\arraybackslash}p{(\linewidth - 12\tabcolsep) * \real{0.1111}}
  >{\raggedright\arraybackslash}p{(\linewidth - 12\tabcolsep) * \real{0.1111}}
  >{\raggedright\arraybackslash}p{(\linewidth - 12\tabcolsep) * \real{0.1111}}@{}}
\toprule\noalign{}
\begin{minipage}[b]{\linewidth}\raggedright
Code
\end{minipage} & \begin{minipage}[b]{\linewidth}\raggedright
Item
\end{minipage} & \begin{minipage}[b]{\linewidth}\raggedright
1
\end{minipage} & \begin{minipage}[b]{\linewidth}\raggedright
2
\end{minipage} & \begin{minipage}[b]{\linewidth}\raggedright
3
\end{minipage} & \begin{minipage}[b]{\linewidth}\raggedright
4
\end{minipage} & \begin{minipage}[b]{\linewidth}\raggedright
5
\end{minipage} \\
\midrule\noalign{}
\endhead
\bottomrule\noalign{}
\endlastfoot
PEOU1 & Les explications SHAP (bar plots, waterfall) sont faciles à
comprendre & 1 & 2 & 3 & 4 & 5 \\
PEOU2 & Le dashboard FRAUDX semble intuitif à prendre en main & 1 & 2 &
3 & 4 & 5 \\
PEOU3 & Un analyste non-spécialiste en IA peut utiliser FRAUDX avec une
formation courte (\textless{} 1 jour) & 1 & 2 & 3 & 4 & 5 \\
PEOU4 & Les visualisations SHAP sont claires même sans background
technique & 1 & 2 & 3 & 4 & 5 \\
PEOU5 & La configuration des seuils d’alerte est accessible & 1 & 2 & 3
& 4 & 5 \\
\end{longtable}
}

Score PEOU = moyenne de PEOU1 à PEOU5 (cible ≥ 3,5/5)

\paragraph{Section D — Confiance (TR)}\label{section-d-confiance-tr}

{\def\LTcaptype{none} % do not increment counter
\begin{longtable}[]{@{}
  >{\raggedright\arraybackslash}p{(\linewidth - 12\tabcolsep) * \real{0.2222}}
  >{\raggedright\arraybackslash}p{(\linewidth - 12\tabcolsep) * \real{0.2222}}
  >{\raggedright\arraybackslash}p{(\linewidth - 12\tabcolsep) * \real{0.1111}}
  >{\raggedright\arraybackslash}p{(\linewidth - 12\tabcolsep) * \real{0.1111}}
  >{\raggedright\arraybackslash}p{(\linewidth - 12\tabcolsep) * \real{0.1111}}
  >{\raggedright\arraybackslash}p{(\linewidth - 12\tabcolsep) * \real{0.1111}}
  >{\raggedright\arraybackslash}p{(\linewidth - 12\tabcolsep) * \real{0.1111}}@{}}
\toprule\noalign{}
\begin{minipage}[b]{\linewidth}\raggedright
Code
\end{minipage} & \begin{minipage}[b]{\linewidth}\raggedright
Item
\end{minipage} & \begin{minipage}[b]{\linewidth}\raggedright
1
\end{minipage} & \begin{minipage}[b]{\linewidth}\raggedright
2
\end{minipage} & \begin{minipage}[b]{\linewidth}\raggedright
3
\end{minipage} & \begin{minipage}[b]{\linewidth}\raggedright
4
\end{minipage} & \begin{minipage}[b]{\linewidth}\raggedright
5
\end{minipage} \\
\midrule\noalign{}
\endhead
\bottomrule\noalign{}
\endlastfoot
TR1 & Je fais confiance aux décisions du modèle ML pour la détection de
fraude & 1 & 2 & 3 & 4 & 5 \\
TR2 & L’explication SHAP renforce ma confiance dans les alertes
produites & 1 & 2 & 3 & 4 & 5 \\
TR3 & Je suis à l’aise avec une détection automatisée sans intervention
humaine préalable & 1 & 2 & 3 & 4 & 5 \\
TR4 & Les mécanismes de sécurité (authentification API, chiffrement,
logs) sont suffisants pour un déploiement bancaire & 1 & 2 & 3 & 4 &
5 \\
TR5 & Je vérifierais moins d’alertes manuellement si le système montrait
une fiabilité constante & 1 & 2 & 3 & 4 & 5 \\
\end{longtable}
}

Score TR = moyenne de TR1 à TR5 (cible ≥ 3,5/5)

\paragraph{Section E — Intention d’adoption
(BI)}\label{section-e-intention-dadoption-bi}

{\def\LTcaptype{none} % do not increment counter
\begin{longtable}[]{@{}
  >{\raggedright\arraybackslash}p{(\linewidth - 12\tabcolsep) * \real{0.2222}}
  >{\raggedright\arraybackslash}p{(\linewidth - 12\tabcolsep) * \real{0.2222}}
  >{\raggedright\arraybackslash}p{(\linewidth - 12\tabcolsep) * \real{0.1111}}
  >{\raggedright\arraybackslash}p{(\linewidth - 12\tabcolsep) * \real{0.1111}}
  >{\raggedright\arraybackslash}p{(\linewidth - 12\tabcolsep) * \real{0.1111}}
  >{\raggedright\arraybackslash}p{(\linewidth - 12\tabcolsep) * \real{0.1111}}
  >{\raggedright\arraybackslash}p{(\linewidth - 12\tabcolsep) * \real{0.1111}}@{}}
\toprule\noalign{}
\begin{minipage}[b]{\linewidth}\raggedright
Code
\end{minipage} & \begin{minipage}[b]{\linewidth}\raggedright
Item
\end{minipage} & \begin{minipage}[b]{\linewidth}\raggedright
1
\end{minipage} & \begin{minipage}[b]{\linewidth}\raggedright
2
\end{minipage} & \begin{minipage}[b]{\linewidth}\raggedright
3
\end{minipage} & \begin{minipage}[b]{\linewidth}\raggedright
4
\end{minipage} & \begin{minipage}[b]{\linewidth}\raggedright
5
\end{minipage} \\
\midrule\noalign{}
\endhead
\bottomrule\noalign{}
\endlastfoot
BI1 & Je recommanderais FRAUDX dans mon institution & 1 & 2 & 3 & 4 &
5 \\
BI2 & Je serais prêt à utiliser FRAUDX quotidiennement & 1 & 2 & 3 & 4 &
5 \\
BI3 & FRAUDX devrait être déployé prioritairement dans les banques
togolaises & 1 & 2 & 3 & 4 & 5 \\
BI4 & Je participerais volontiers à une phase pilote & 1 & 2 & 3 & 4 &
5 \\
BI5 & Je soutiendrais l’allocation d’un budget pour ce système & 1 & 2 &
3 & 4 & 5 \\
\end{longtable}
}

Score BI = moyenne de BI1 à BI5 (cible ≥ 3,5/5)

\paragraph{Section F — Facteurs contextuels
(FC)}\label{section-f-facteurs-contextuels-fc}

{\def\LTcaptype{none} % do not increment counter
\begin{longtable}[]{@{}
  >{\raggedright\arraybackslash}p{(\linewidth - 12\tabcolsep) * \real{0.2222}}
  >{\raggedright\arraybackslash}p{(\linewidth - 12\tabcolsep) * \real{0.2222}}
  >{\raggedright\arraybackslash}p{(\linewidth - 12\tabcolsep) * \real{0.1111}}
  >{\raggedright\arraybackslash}p{(\linewidth - 12\tabcolsep) * \real{0.1111}}
  >{\raggedright\arraybackslash}p{(\linewidth - 12\tabcolsep) * \real{0.1111}}
  >{\raggedright\arraybackslash}p{(\linewidth - 12\tabcolsep) * \real{0.1111}}
  >{\raggedright\arraybackslash}p{(\linewidth - 12\tabcolsep) * \real{0.1111}}@{}}
\toprule\noalign{}
\begin{minipage}[b]{\linewidth}\raggedright
Code
\end{minipage} & \begin{minipage}[b]{\linewidth}\raggedright
Item
\end{minipage} & \begin{minipage}[b]{\linewidth}\raggedright
1
\end{minipage} & \begin{minipage}[b]{\linewidth}\raggedright
2
\end{minipage} & \begin{minipage}[b]{\linewidth}\raggedright
3
\end{minipage} & \begin{minipage}[b]{\linewidth}\raggedright
4
\end{minipage} & \begin{minipage}[b]{\linewidth}\raggedright
5
\end{minipage} \\
\midrule\noalign{}
\endhead
\bottomrule\noalign{}
\endlastfoot
FC1 & Les contraintes infrastructurelles (électricité, bande passante)
sont un frein majeur & 1 & 2 & 3 & 4 & 5 \\
FC2 & Le manque de données locales de fraude est un obstacle important &
1 & 2 & 3 & 4 & 5 \\
FC3 & La conformité réglementaire (BCEAO, CNIL) est une préoccupation
prioritaire & 1 & 2 & 3 & 4 & 5 \\
FC4 & Mon institution dispose des compétences techniques nécessaires
pour exploiter un système IA & 1 & 2 & 3 & 4 & 5 \\
FC5 & La direction est généralement ouverte à l’innovation technologique
& 1 & 2 & 3 & 4 & 5 \\
\end{longtable}
}

Score FC = items FC1-FC3 inversés (1→5, 5→1) puis moyenne avec FC4-FC5

\subsubsection{Partie 4 — Grille d’analyse des résultats
TAM}\label{partie-4-grille-danalyse-des-ruxe9sultats-tam}

\paragraph{4.1. Grille de codification des entretiens
semi-structurés}\label{grille-de-codification-des-entretiens-semi-structuruxe9s}

{\def\LTcaptype{none} % do not increment counter
\begin{longtable}[]{@{}
  >{\raggedright\arraybackslash}p{(\linewidth - 8\tabcolsep) * \real{0.1034}}
  >{\raggedright\arraybackslash}p{(\linewidth - 8\tabcolsep) * \real{0.1897}}
  >{\raggedright\arraybackslash}p{(\linewidth - 8\tabcolsep) * \real{0.1207}}
  >{\raggedright\arraybackslash}p{(\linewidth - 8\tabcolsep) * \real{0.3621}}
  >{\raggedright\arraybackslash}p{(\linewidth - 8\tabcolsep) * \real{0.2241}}@{}}
\toprule\noalign{}
\begin{minipage}[b]{\linewidth}\raggedright
Code
\end{minipage} & \begin{minipage}[b]{\linewidth}\raggedright
Dimension
\end{minipage} & \begin{minipage}[b]{\linewidth}\raggedright
Thème
\end{minipage} & \begin{minipage}[b]{\linewidth}\raggedright
Mots-clés / Indices
\end{minipage} & \begin{minipage}[b]{\linewidth}\raggedright
Pondération
\end{minipage} \\
\midrule\noalign{}
\endhead
\bottomrule\noalign{}
\endlastfoot
PU+ & Utilité perçue — positive & Thème 2 & « améliorerait », « utile »,
« gain de temps », « nécessaire » & +1 \\
PU- & Utilité perçue — négative & Thème 2 & « pas convaincu », « déjà
suffisant », « inutile » & -1 \\
PEOU+ & Facilité — positive & Thème 3 & « simple », « intuitif », «
clair », « rapide à apprendre » & +1 \\
PEOU- & Facilité — négative & Thème 3 & « complexe », « trop technique
», « besoin de formation longue » & -1 \\
TR+ & Confiance — positive & Thème 4 & « fiable », « confiance », «
transparent », « explique bien » & +1 \\
TR- & Confiance — négative & Thème 4 & « boîte noire », « méfiance », «
pas fiable », « risques » & -1 \\
BI+ & Adoption — positive & Thème 5 & « oui », « je recommande », « prêt
», « prioritaire » & +1 \\
BI- & Adoption — négative & Thème 5 & « pas maintenant », « trop tôt »,
« pas prioritaire » & -1 \\
FC+ & Contexte — facilitateur & Thème 5 & « infrastructure OK », «
compétences disponibles », « budget » & +1 \\
FC- & Contexte — frein & Thème 5 & « pas d’électricité », « problème de
données », « coût » & -1 \\
\end{longtable}
}

\paragraph{4.2. Scoring des entretiens}\label{scoring-des-entretiens}

{\def\LTcaptype{none} % do not increment counter
\begin{longtable}[]{@{}
  >{\raggedright\arraybackslash}p{(\linewidth - 4\tabcolsep) * \real{0.2549}}
  >{\raggedright\arraybackslash}p{(\linewidth - 4\tabcolsep) * \real{0.4314}}
  >{\raggedright\arraybackslash}p{(\linewidth - 4\tabcolsep) * \real{0.3137}}@{}}
\toprule\noalign{}
\begin{minipage}[b]{\linewidth}\raggedright
Score total
\end{minipage} & \begin{minipage}[b]{\linewidth}\raggedright
Niveau d’acceptabilité
\end{minipage} & \begin{minipage}[b]{\linewidth}\raggedright
Interprétation
\end{minipage} \\
\midrule\noalign{}
\endhead
\bottomrule\noalign{}
\endlastfoot
≥ +8 & Élevé & FRAUDX bien accepté, conditions réunies pour un pilote \\
+3 à +7 & Modéré & Acceptable sous réserves (formation, adaptation) \\
-2 à +2 & Neutre & Attentiste, besoin de démonstration complémentaire \\
≤ -3 & Faible & Freins majeurs, revoir la stratégie de déploiement \\
\end{longtable}
}

\paragraph{4.3. Analyse croisée quantitative +
qualitative}\label{analyse-croisuxe9e-quantitative-qualitative}

La validation de HS3 combine les deux volets :

{\def\LTcaptype{none} % do not increment counter
\begin{longtable}[]{@{}
  >{\raggedright\arraybackslash}p{(\linewidth - 4\tabcolsep) * \real{0.1892}}
  >{\raggedright\arraybackslash}p{(\linewidth - 4\tabcolsep) * \real{0.2162}}
  >{\raggedright\arraybackslash}p{(\linewidth - 4\tabcolsep) * \real{0.5946}}@{}}
\toprule\noalign{}
\begin{minipage}[b]{\linewidth}\raggedright
Seuil
\end{minipage} & \begin{minipage}[b]{\linewidth}\raggedright
Source
\end{minipage} & \begin{minipage}[b]{\linewidth}\raggedright
Statut dans ce mémoire
\end{minipage} \\
\midrule\noalign{}
\endhead
\bottomrule\noalign{}
\endlastfoot
FP ≤ 2\% & Modèle XGBoost (quantitatif) & 4,77\% — partiellement
atteint \\
Satisfaction PU ≥ 70\% (≥ 3,5/5) & Questionnaire TAM (section B) &
Perspective — non administré \\
Majorité de codes PU+/PEOU+/TR+ & Entretiens semi-structurés (grille
4.1) & Perspective — non administré \\
\end{longtable}
}

La convergence entre score TAM ≥ 3,5/5, codification entretien ≥ +3 et
FP ≤ 2\% constituerait une validation empirique forte de HS3.

\begin{center}\rule{0.5\linewidth}{0.5pt}\end{center}

\subsection{Annexe C : Description détaillée des indicateurs et formules
de
calcul}\label{annexe-c-description-duxe9tailluxe9e-des-indicateurs-et-formules-de-calcul}

\subsubsection{C.1. Variables transactionnelles et
temporelles}\label{c.1.-variables-transactionnelles-et-temporelles}

{\def\LTcaptype{none} % do not increment counter
\begin{longtable}[]{@{}
  >{\raggedright\arraybackslash}p{(\linewidth - 6\tabcolsep) * \real{0.2128}}
  >{\raggedright\arraybackslash}p{(\linewidth - 6\tabcolsep) * \real{0.1277}}
  >{\raggedright\arraybackslash}p{(\linewidth - 6\tabcolsep) * \real{0.4894}}
  >{\raggedright\arraybackslash}p{(\linewidth - 6\tabcolsep) * \real{0.1702}}@{}}
\toprule\noalign{}
\begin{minipage}[b]{\linewidth}\raggedright
Variable
\end{minipage} & \begin{minipage}[b]{\linewidth}\raggedright
Type
\end{minipage} & \begin{minipage}[b]{\linewidth}\raggedright
Formule / Description
\end{minipage} & \begin{minipage}[b]{\linewidth}\raggedright
Source
\end{minipage} \\
\midrule\noalign{}
\endhead
\bottomrule\noalign{}
\endlastfoot
\texttt{TransactionAmt} & Continue & Montant brut de la transaction en
USD & IEEE-CIS \\
\texttt{TransactionDT} & Temporelle & Timestamp brut (secondes écoulées
depuis le 01/12/2017) & IEEE-CIS \\
\texttt{hour} & Catégorielle &
\texttt{hour\ =\ (TransactionDT\ //\ 3600)\ \%\ 24} & Feature
engineering \\
\texttt{dayofweek} & Catégorielle &
\texttt{dayofweek\ =\ (TransactionDT\ //\ 86400)\ \%\ 7} & Feature
engineering \\
\texttt{is\_night} & Binaire &
\texttt{1\ if\ (hour\ \textless{}\ 6\ or\ hour\ \textgreater{}\ 22)\ else\ 0}
& Feature engineering \\
\texttt{is\_weekend} & Binaire &
\texttt{1\ if\ dayofweek\ in\ \{5,\ 6\}\ else\ 0} & Feature
engineering \\
\texttt{ProductCD} & Catégorielle & Type de produit (C = carte, W =
wallet, etc.) & IEEE-CIS \\
\end{longtable}
}

\subsubsection{C.2. Variables de carte et
dispositif}\label{c.2.-variables-de-carte-et-dispositif}

{\def\LTcaptype{none} % do not increment counter
\begin{longtable}[]{@{}
  >{\raggedright\arraybackslash}p{(\linewidth - 4\tabcolsep) * \real{0.3448}}
  >{\raggedright\arraybackslash}p{(\linewidth - 4\tabcolsep) * \real{0.2069}}
  >{\raggedright\arraybackslash}p{(\linewidth - 4\tabcolsep) * \real{0.4483}}@{}}
\toprule\noalign{}
\begin{minipage}[b]{\linewidth}\raggedright
Variable
\end{minipage} & \begin{minipage}[b]{\linewidth}\raggedright
Type
\end{minipage} & \begin{minipage}[b]{\linewidth}\raggedright
Description
\end{minipage} \\
\midrule\noalign{}
\endhead
\bottomrule\noalign{}
\endlastfoot
\texttt{card1} à \texttt{card6} & Catégorielle & Identifiants anonymisés
de la carte (émetteur, type, catégorie) \\
\texttt{addr1}, \texttt{addr2} & Catégorielle & Codes de localisation
anonymisés (pays, région) \\
\texttt{dist1} & Continue & Distance anonymisée entre adresse de
facturation et IP \\
\texttt{P\_emaildomain}, \texttt{R\_emaildomain} & Catégorielle &
Domaines email de l’acheteur et du destinataire \\
\end{longtable}
}

\subsubsection{C.3. Variables de comportement et
vélocité}\label{c.3.-variables-de-comportement-et-vuxe9locituxe9}

{\def\LTcaptype{none} % do not increment counter
\begin{longtable}[]{@{}
  >{\raggedright\arraybackslash}p{(\linewidth - 4\tabcolsep) * \real{0.2564}}
  >{\raggedright\arraybackslash}p{(\linewidth - 4\tabcolsep) * \real{0.1538}}
  >{\raggedright\arraybackslash}p{(\linewidth - 4\tabcolsep) * \real{0.5897}}@{}}
\toprule\noalign{}
\begin{minipage}[b]{\linewidth}\raggedright
Variable
\end{minipage} & \begin{minipage}[b]{\linewidth}\raggedright
Type
\end{minipage} & \begin{minipage}[b]{\linewidth}\raggedright
Formule / Description
\end{minipage} \\
\midrule\noalign{}
\endhead
\bottomrule\noalign{}
\endlastfoot
\texttt{tx\_count\_by\_card1} & Compteur &
\texttt{count(transactions\ group\ by\ card1)} sur intervalle
glissant \\
\texttt{avg\_amount\_by\_card1} & Continue &
\texttt{mean(TransactionAmt\ group\ by\ card1)} sur historique \\
\texttt{log\_amount} & Continue & \texttt{log(TransactionAmt\ +\ 1)} —
normalisation logarithmique \\
\texttt{amt\_diff\_from\_avg} & Continue &
\texttt{TransactionAmt\ -\ avg\_amount\_by\_card1} — écart au profil \\
C* (C1 à C14) & Continue & Variables calculées par l’émetteur (scores de
risque) \\
D* (D1 à D15) & Continue & Délais anonymisés entre événements
(secondes) \\
\end{longtable}
}

\subsubsection{C.4. Variables anonymisées
(PCA)}\label{c.4.-variables-anonymisuxe9es-pca}

Les variables V* (V1 à V339) sont issues d’une transformation PCA
(analyse en composantes principales) appliquée par les émetteurs pour
anonymiser des données sensibles sans perte d’information discriminante.
Ces variables ne sont pas interprétables individuellement mais leur
contribution aux valeurs SHAP peut être analysée globalement.

\begin{center}\rule{0.5\linewidth}{0.5pt}\end{center}

\subsection{Annexe D : Questions anticipées de
soutenance}\label{annexe-d-questions-anticipuxe9es-de-soutenance}

Cette annexe recense les questions probables que le jury pourrait poser
lors de la soutenance, avec des éléments de réponse structurés renvoyant
aux sections correspondantes du mémoire.

\subsubsection{D.1. Questions sur le cadrage et la
problématique}\label{d.1.-questions-sur-le-cadrage-et-la-probluxe9matique}

\textbf{Q1 — Pourquoi avoir choisi le jeu de données IEEE-CIS plutôt que
des données togolaises ?}

Faute de données bancaires réelles togolaises (confidentialité, absence
de partenariat avec une institution financière locale), IEEE-CIS
constitue le proxy le plus proche disponible. Ce dataset présente des
caractéristiques transactionnelles transférables au Togo (montants,
temporalité, vélocité). Les limites de cette approximation sont
explicitement discutées en III.1.4.

\textbf{Q2 — En quoi votre approche est-elle originale par rapport à la
littérature existante ?}

Trois originalités : (1) première étude documentée sur le Togo combinant
banque et mobile money ; (2) architecture complète allant du modèle ML
au dashboard sécurisé (JWT + RBAC) ; (3) validation empirique de
l’explicabilité SHAP comme levier de réduction des FP. Voir I.3 et la
synthèse du Tableau I.1.

\subsubsection{D.2. Questions sur la
méthodologie}\label{d.2.-questions-sur-la-muxe9thodologie}

\textbf{Q3 — Pourquoi trois modèles seulement ? Pourquoi pas de Deep
Learning ?}

Le choix de trois modèles complémentaires (non supervisé, supervisé
ensembliste, supervisé boosting) couvre les paradigmes principaux. Le
Deep Learning (LSTM, Transformers) a été exclu pour trois raisons :
absence de GPU, absence de structure séquentielle explicite dans les
données, objectif de déploiement sur infrastructure modeste. Cette
limitation est discutée en III.4.1.

\textbf{Q4 — Pourquoi SMOTE plutôt que d’autres techniques de
rééquilibrage (ADASYN, Random Under-sampling, cost-sensitive learning)
?}

SMOTE a été retenu pour son efficacité éprouvée sur les données
tabulaires déséquilibrées et sa disponibilité dans Imbalanced-learn.
ADASYN, plus sensible au bruit, n’a pas montré de supériorité
significative dans la littérature sur des données similaires.
L’under-sampling aléatoire aurait perdu trop d’exemples de la classe
majoritaire. Le cost-sensitive learning (via \texttt{scale\_pos\_weight}
dans XGBoost) a été combiné à SMOTE, cumulant les avantages des deux
approches. Voir II.4.

\textbf{Q5 — Pourquoi avoir utilisé un holdout simple plutôt qu’une
validation croisée (k-fold) ?}

Le holdout simple 80/20 avec stratification a été préféré à la
validation croisée pour trois raisons : (1) la taille du jeu de données
(118 108 transactions de test) garantit une estimation fiable sans
nécessiter de ré-échantillonnage ; (2) Optuna utilise une validation
croisée 3-folds sur le seul ensemble d’entraînement, combinant les
avantages des deux approches ; (3) le holdout permet une interprétation
directe de la matrice de confusion sur un ensemble de test fixe. Voir
III.2.1.

\subsubsection{D.3. Questions sur les
résultats}\label{d.3.-questions-sur-les-ruxe9sultats}

\textbf{Q6 — Le F1-Score de 0,23 est très faible. Comment justifiez-vous
la performance du modèle ?}

Cette question est centrale. Le F1 de 0,23 reflète l’arbitrage délibéré
en faveur du Recall (85,02 \%) au détriment de la Précision (13,54 \%).
Dans le contexte de la détection de fraude bancaire, un faux négatif
(fraude non détectée) coûte en moyenne 10 à 100 fois plus qu’un faux
positif (alerte superflue). Le système FRAUDX est conçu comme un outil
de priorisation d’alertes pour analystes, non comme un système de
blocage automatique. Les 20,7 \% de FP génèrent un volume d’alertes
gérable par une équipe de 3 à 5 analystes, chaque alerte étant
accompagnée d’une explication SHAP permettant une qualification en moins
de 30 secondes. Voir III.4.1 et IV.6.2.

\textbf{Q7 — Pourquoi Random Forest a-t-il un meilleur F1 (0,44) que
XGBoost (0,23) mais un Recall inférieur (0,62 contre 0,85) ?}

Random Forest optimise naturellement l’équilibre précision/rappel via le
vote majoritaire de ses arbres, ce qui produit un F1 plus élevé mais au
prix d’un nombre important de faux négatifs (38 \% de fraudes non
détectées). XGBoost, avec \texttt{scale\_pos\_weight=22,4} et un seuil
abaissé à 0,35, sacrifie la précision pour maximiser le Recall. Le choix
entre les deux dépend du contexte opérationnel : pour une détection
maximale, XGBoost est préféré ; pour un équilibre, Random Forest
pourrait être retenu. Voir Tableau III.1.

\textbf{Q8 — L’AUC-PR de 0,57 est inférieure à la cible de 0,65.
Qu’est-ce que cela implique ?}

L’AUC-PR mesure la performance moyenne du modèle sur tous les seuils
possibles. Un score de 0,57 indique que le modèle discrimine
correctement une partie significative des transactions frauduleuses mais
que la marge de séparation entre classes reste limitée en raison du fort
déséquilibre et de l’absence de variables spécifiques au contexte
togolais. L’optimisation du seuil à 0,35 améliore le Recall mais ne peut
compenser structurellement la qualité intrinsèque du scoring. L’AUC-PR
pourrait être améliorée par l’ajout de données locales. Voir III.4.1.

\subsubsection{D.4. Questions sur la plateforme et
l’architecture}\label{d.4.-questions-sur-la-plateforme-et-larchitecture}

\textbf{Q9 — Comment garantissez-vous la sécurité des données bancaires
?}

Trois niveaux de sécurité : (1) authentification JWT (HMAC-SHA256) avec
tokens limités dans le temps (30 minutes) ; (2) contrôle d’accès RBAC
définissant strictement les permissions par rôle (analyste, superviseur,
administrateur) ; (3) chiffrement AES-256 pour les données au repos et
TLS 1.3 pour les données en transit. Un limiteur de débit (100
req/min/IP) protège contre les abus. La conformité à la loi N°2020-003
et aux directives BCEAO est assurée par des logs d’audit horodatés
conservés 10 ans. Voir III.3.3.

\textbf{Q10 — Le prototype est-il industrialisable ? Quelles sont les
étapes nécessaires ?}

Le prototype démontre la faisabilité technique mais n’est pas
industrialisable en l’état. Les étapes nécessaires sont : (1)
conteneurisation (Docker, Kubernetes) et déploiement cloud/hybride ; (2)
migration vers une base de données relationnelle robuste (PostgreSQL) ;
(3) mise en place d’un pipeline CI/CD ; (4) tests de charge (cible : 50
000 transactions/jour) ; (5) ajout d’un WAF applicatif ; (6)
certification par un auditeur externe. Un budget de 72 500 € est estimé
pour la première année (IV.6.1).

\subsubsection{D.5. Questions sur les
perspectives}\label{d.5.-questions-sur-les-perspectives}

\textbf{Q11 — Comment passer de 20,7 \% de FP à la cible de 2 \% ?}

Plusieurs leviers : (1) réentraînement avec feedback des analystes
(\textasciitilde3 mois pour observer une réduction significative) ; (2)
ajout de données locales togolaises pour enrichir la représentation des
transactions légitimes ; (3) calibrage plus fin du seuil par transaction
(seuil dynamique selon le montant, l’historique client) ; (4) ajout
d’une couche de règles métier post-ML pour filtrer les FP évidents. Voir
IV.4.6.

\textbf{Q12 — Quelles sont les prochaines étapes de votre recherche ?}

À court terme : établir un partenariat avec une banque togolaise pour
valider le modèle sur des données réelles. À moyen terme : déployer le
pilote FRAUDX dans une institution partenaire (SUNU Bank ou autre) et
mesurer l’acceptabilité via le questionnaire TAM (Annexe B). À long
terme : extension à l’échelle de l’UEMOA et intégration de techniques
avancées (apprentissage fédéré, LSTM). Voir IV.7.

\begin{center}\rule{0.5\linewidth}{0.5pt}\end{center}

\section{TABLE DES MATIERES}\label{table-des-matieres}

DEDICACE ……………………………………………………………………………………………………………. I REMERCIEMENTS
……………………………………………………………………………………………… II RESUME
…………………………………………………………………………………………………………….. III ABSTRACT
………………………………………………………………………………………………………… IV SOMMAIRE
…………………………………………………………………………………………………………. V LISTE DES TABLEAUX
……………………………………………………………………………………… VI LISTE DES FIGURES ET GRAPHIQUES
…………………………………………………………….. VII LISTE DES ABREVIATIONS
……………………………………………………………………………. VIII

INTRODUCTION GENERALE ………………………………………………………………………………. 1 1. CONTEXTE
GENERAL DE L’ETUDE ……………………………………………………………….. 2 2. PROBLEMATIQUE DE
L’ETUDE …………………………………………………………………….. 3 2.1. Présentation du problème
…………………………………………………………………………… 3 2.2. Formulation du problème
…………………………………………………………………………… 3 3. HYPOTHESES DE L’ETUDE
……………………………………………………………………………. 4 3.1. Hypothèse générale
……………………………………………………………………………………. 4 3.2. Hypothèses spécifiques
………………………………………………………………………………. 4 4. OBJECTIFS DE L’ETUDE
…………………………………………………………………………………. 5 4.1. Objectif général
…………………………………………………………………………………………. 5 4.2. Objectifs spécifiques
………………………………………………………………………………….. 5 5. JUSTIFICATION DE L’ETUDE
………………………………………………………………………….. 6 6. DELIMITATION DE L’ETUDE
………………………………………………………………………….. 7 7. PLAN DU MEMOIRE
………………………………………………………………………………………. 8

CHAPITRE I : CADRE THEORIQUE ET CONCEPTUEL ………………………………………… 9
Introduction ……………………………………………………………………………………………………… 10 I.1. Cadre
théorique et état de l’art …………………………………………………………………….. 11 I.1.1. La
fraude bancaire et mobile money ………………………………………………………. 11 I.1.2. Le
Machine Learning appliqué à la détection de fraude ……………………………. 14
I.1.3. L’explicabilité (XAI) ……………………………………………………………………………. 18 I.1.4.
Cadre légal et réglementaire …………………………………………………………………. 20 I.2.
Historique et évolution du domaine ………………………………………………………………. 22 I.3.
Études antérieures et lacunes ………………………………………………………………………. 24 Conclusion
du chapitre ……………………………………………………………………………………… 27

CHAPITRE II : METHODOLOGIE DE L’ETUDE ………………………………………………….. 28
Introduction ……………………………………………………………………………………………………… 28 II.1. Nature de
l’étude ………………………………………………………………………………………. 29 II.2. Variables de l’étude
…………………………………………………………………………………… 30 II.2.4. Dynamique anticipée des
variables …………………………………………………………. 31 II.3. Population et échantillon
……………………………………………………………………………. 32 II.4. Approche méthodologique
………………………………………………………………………….. 33 II.5. Outils de l’étude
………………………………………………………………………………………. 35 II.6. Stratégie de vérification
des hypothèses ……………………………………………………….. 36 Conclusion du chapitre
……………………………………………………………………………………… 36

CHAPITRE III : PRESENTATION DE LA SITUATION ET COLLECTE DES DONNEES ..
37 Introduction ……………………………………………………………………………………………………… 37 III.1.
Présentation et analyse exploratoire des données …………………………………………. 38
III.2. Conception et évaluation des modèles ………………………………………………………… 40
III.3. Proposition de plateforme FRAUDX …………………………………………………………… 43
III.4. Tests et validation ……………………………………………………………………………………. 46
Conclusion du chapitre ……………………………………………………………………………………… 48

CHAPITRE IV : ANALYSE-DIAGNOSTIC ET PROPOSITION D’INTERVENTION ……. 49
Introduction ……………………………………………………………………………………………………… 49 IV.1.
Présentation et analyse de la situation ………………………………………………………….. 50
IV.1.1. Méthodologie d’investigation ……………………………………………………………….. 50
IV.1.2. Diagnostic de la situation …………………………………………………………………….. 50
IV.1.3. Synthèse diagnostique et besoins stratégiques ……………………………………………
51 IV.1.4. Vérification des hypothèses ………………………………………………………………….. 52
IV.2. Intervention proposée et justification ………………………………………………………….. 53
IV.2.1. L’intervention : le système FRAUDX ………………………………………………………. 53
IV.2.2. Justification de l’intervention ………………………………………………………………. 53
IV.3. Objectifs de l’intervention …………………………………………………………………………. 54
IV.3.1. Objectif général ………………………………………………………………………………… 54 IV.3.2.
Objectifs spécifiques ………………………………………………………………………….. 54 IV.4.
Composantes de l’intervention envisagée ……………………………………………………….. 55
IV.4.1. Module de collecte et prétraitement ……………………………………………………….. 55
IV.4.2. Moteur de détection XGBoost ……………………………………………………………….. 55
IV.4.3. Module d’explicabilité SHAP ………………………………………………………………… 56 IV.4.4.
Dashboard interactif ……………………………………………………………………….. 56 IV.4.5. Module de
sécurité et conformité …………………………………………………………… 56 IV.4.6. Module de
feedback et apprentissage continu …………………………………………….. 57 IV.4.7. API
d’intégration et connecteurs …………………………………………………………… 57 IV.5. Stratégies
d’action et périmètre …………………………………………………………………. 58 IV.5.1. Stratégies
d’action …………………………………………………………………………… 58 IV.5.2. Périmètre de
l’intervention …………………………………………………………………. 58 IV.6. Étude de faisabilité
…………………………………………………………………………………. 59 IV.6.1. Faisabilité économique
……………………………………………………………………….. 59 IV.6.2. Faisabilité sociale
……………………………………………………………………………. 60 IV.6.3. Faisabilité technique
…………………………………………………………………………. 60 IV.6.4. Faisabilité environnementale
………………………………………………………………… 61 IV.7. Perspectives et limites
……………………………………………………………………………. 62 Conclusion du chapitre
……………………………………………………………………………………… 62

CONCLUSION GENERALE ………………………………………………………………………………… 64 BIBLIOGRAPHIE ET
WEBOGRAPHIE ………………………………………………………………….. X ANNEXES
………………………………………………………………………………………………………… XIII Annexe A : Grille
d’entretien semi-directif ……………………………………………………… XIII Annexe B : Guide
TAM complet — Adoption du système FRAUDX …………………… XVI Annexe C :
Description détaillée des indicateurs et formules de calcul ………………..
XXII Annexe D : Questions anticipées de soutenance …………………………………………………
XXV

\end{document}
